% Options for packages loaded elsewhere
\PassOptionsToPackage{unicode}{hyperref}
\PassOptionsToPackage{hyphens}{url}
\PassOptionsToPackage{dvipsnames,svgnames,x11names}{xcolor}
%
\documentclass[
  11pt,
]{article}

\usepackage{amsmath,amssymb}
\usepackage{iftex}
\ifPDFTeX
  \usepackage[T1]{fontenc}
  \usepackage[utf8]{inputenc}
  \usepackage{textcomp} % provide euro and other symbols
\else % if luatex or xetex
  \usepackage{unicode-math}
  \defaultfontfeatures{Scale=MatchLowercase}
  \defaultfontfeatures[\rmfamily]{Ligatures=TeX,Scale=1}
\fi
\usepackage{lmodern}
\ifPDFTeX\else  
    % xetex/luatex font selection
    \setmainfont[]{Arial}
\fi
% Use upquote if available, for straight quotes in verbatim environments
\IfFileExists{upquote.sty}{\usepackage{upquote}}{}
\IfFileExists{microtype.sty}{% use microtype if available
  \usepackage[]{microtype}
  \UseMicrotypeSet[protrusion]{basicmath} % disable protrusion for tt fonts
}{}
\makeatletter
\@ifundefined{KOMAClassName}{% if non-KOMA class
  \IfFileExists{parskip.sty}{%
    \usepackage{parskip}
  }{% else
    \setlength{\parindent}{0pt}
    \setlength{\parskip}{6pt plus 2pt minus 1pt}}
}{% if KOMA class
  \KOMAoptions{parskip=half}}
\makeatother
\usepackage{xcolor}
\usepackage[a4paper, margin=0.8in]{geometry}
\setlength{\emergencystretch}{3em} % prevent overfull lines
\setcounter{secnumdepth}{5}
% Make \paragraph and \subparagraph free-standing
\makeatletter
\ifx\paragraph\undefined\else
  \let\oldparagraph\paragraph
  \renewcommand{\paragraph}{
    \@ifstar
      \xxxParagraphStar
      \xxxParagraphNoStar
  }
  \newcommand{\xxxParagraphStar}[1]{\oldparagraph*{#1}\mbox{}}
  \newcommand{\xxxParagraphNoStar}[1]{\oldparagraph{#1}\mbox{}}
\fi
\ifx\subparagraph\undefined\else
  \let\oldsubparagraph\subparagraph
  \renewcommand{\subparagraph}{
    \@ifstar
      \xxxSubParagraphStar
      \xxxSubParagraphNoStar
  }
  \newcommand{\xxxSubParagraphStar}[1]{\oldsubparagraph*{#1}\mbox{}}
  \newcommand{\xxxSubParagraphNoStar}[1]{\oldsubparagraph{#1}\mbox{}}
\fi
\makeatother

\usepackage{color}
\usepackage{fancyvrb}
\newcommand{\VerbBar}{|}
\newcommand{\VERB}{\Verb[commandchars=\\\{\}]}
\DefineVerbatimEnvironment{Highlighting}{Verbatim}{commandchars=\\\{\}}
% Add ',fontsize=\small' for more characters per line
\usepackage{framed}
\definecolor{shadecolor}{RGB}{248,248,248}
\newenvironment{Shaded}{\begin{snugshade}}{\end{snugshade}}
\newcommand{\AlertTok}[1]{\textcolor[rgb]{0.94,0.16,0.16}{#1}}
\newcommand{\AnnotationTok}[1]{\textcolor[rgb]{0.56,0.35,0.01}{\textbf{\textit{#1}}}}
\newcommand{\AttributeTok}[1]{\textcolor[rgb]{0.13,0.29,0.53}{#1}}
\newcommand{\BaseNTok}[1]{\textcolor[rgb]{0.00,0.00,0.81}{#1}}
\newcommand{\BuiltInTok}[1]{#1}
\newcommand{\CharTok}[1]{\textcolor[rgb]{0.31,0.60,0.02}{#1}}
\newcommand{\CommentTok}[1]{\textcolor[rgb]{0.56,0.35,0.01}{\textit{#1}}}
\newcommand{\CommentVarTok}[1]{\textcolor[rgb]{0.56,0.35,0.01}{\textbf{\textit{#1}}}}
\newcommand{\ConstantTok}[1]{\textcolor[rgb]{0.56,0.35,0.01}{#1}}
\newcommand{\ControlFlowTok}[1]{\textcolor[rgb]{0.13,0.29,0.53}{\textbf{#1}}}
\newcommand{\DataTypeTok}[1]{\textcolor[rgb]{0.13,0.29,0.53}{#1}}
\newcommand{\DecValTok}[1]{\textcolor[rgb]{0.00,0.00,0.81}{#1}}
\newcommand{\DocumentationTok}[1]{\textcolor[rgb]{0.56,0.35,0.01}{\textbf{\textit{#1}}}}
\newcommand{\ErrorTok}[1]{\textcolor[rgb]{0.64,0.00,0.00}{\textbf{#1}}}
\newcommand{\ExtensionTok}[1]{#1}
\newcommand{\FloatTok}[1]{\textcolor[rgb]{0.00,0.00,0.81}{#1}}
\newcommand{\FunctionTok}[1]{\textcolor[rgb]{0.13,0.29,0.53}{\textbf{#1}}}
\newcommand{\ImportTok}[1]{#1}
\newcommand{\InformationTok}[1]{\textcolor[rgb]{0.56,0.35,0.01}{\textbf{\textit{#1}}}}
\newcommand{\KeywordTok}[1]{\textcolor[rgb]{0.13,0.29,0.53}{\textbf{#1}}}
\newcommand{\NormalTok}[1]{#1}
\newcommand{\OperatorTok}[1]{\textcolor[rgb]{0.81,0.36,0.00}{\textbf{#1}}}
\newcommand{\OtherTok}[1]{\textcolor[rgb]{0.56,0.35,0.01}{#1}}
\newcommand{\PreprocessorTok}[1]{\textcolor[rgb]{0.56,0.35,0.01}{\textit{#1}}}
\newcommand{\RegionMarkerTok}[1]{#1}
\newcommand{\SpecialCharTok}[1]{\textcolor[rgb]{0.81,0.36,0.00}{\textbf{#1}}}
\newcommand{\SpecialStringTok}[1]{\textcolor[rgb]{0.31,0.60,0.02}{#1}}
\newcommand{\StringTok}[1]{\textcolor[rgb]{0.31,0.60,0.02}{#1}}
\newcommand{\VariableTok}[1]{\textcolor[rgb]{0.00,0.00,0.00}{#1}}
\newcommand{\VerbatimStringTok}[1]{\textcolor[rgb]{0.31,0.60,0.02}{#1}}
\newcommand{\WarningTok}[1]{\textcolor[rgb]{0.56,0.35,0.01}{\textbf{\textit{#1}}}}

\providecommand{\tightlist}{%
  \setlength{\itemsep}{0pt}\setlength{\parskip}{0pt}}\usepackage{longtable,booktabs,array}
\usepackage{calc} % for calculating minipage widths
% Correct order of tables after \paragraph or \subparagraph
\usepackage{etoolbox}
\makeatletter
\patchcmd\longtable{\par}{\if@noskipsec\mbox{}\fi\par}{}{}
\makeatother
% Allow footnotes in longtable head/foot
\IfFileExists{footnotehyper.sty}{\usepackage{footnotehyper}}{\usepackage{footnote}}
\makesavenoteenv{longtable}
\usepackage{graphicx}
\makeatletter
\newsavebox\pandoc@box
\newcommand*\pandocbounded[1]{% scales image to fit in text height/width
  \sbox\pandoc@box{#1}%
  \Gscale@div\@tempa{\textheight}{\dimexpr\ht\pandoc@box+\dp\pandoc@box\relax}%
  \Gscale@div\@tempb{\linewidth}{\wd\pandoc@box}%
  \ifdim\@tempb\p@<\@tempa\p@\let\@tempa\@tempb\fi% select the smaller of both
  \ifdim\@tempa\p@<\p@\scalebox{\@tempa}{\usebox\pandoc@box}%
  \else\usebox{\pandoc@box}%
  \fi%
}
% Set default figure placement to htbp
\def\fps@figure{htbp}
\makeatother

\makeatletter
\@ifpackageloaded{caption}{}{\usepackage{caption}}
\AtBeginDocument{%
\ifdefined\contentsname
  \renewcommand*\contentsname{Table of contents}
\else
  \newcommand\contentsname{Table of contents}
\fi
\ifdefined\listfigurename
  \renewcommand*\listfigurename{List of Figures}
\else
  \newcommand\listfigurename{List of Figures}
\fi
\ifdefined\listtablename
  \renewcommand*\listtablename{List of Tables}
\else
  \newcommand\listtablename{List of Tables}
\fi
\ifdefined\figurename
  \renewcommand*\figurename{Figure}
\else
  \newcommand\figurename{Figure}
\fi
\ifdefined\tablename
  \renewcommand*\tablename{Table}
\else
  \newcommand\tablename{Table}
\fi
}
\@ifpackageloaded{float}{}{\usepackage{float}}
\floatstyle{ruled}
\@ifundefined{c@chapter}{\newfloat{codelisting}{h}{lop}}{\newfloat{codelisting}{h}{lop}[chapter]}
\floatname{codelisting}{Listing}
\newcommand*\listoflistings{\listof{codelisting}{List of Listings}}
\makeatother
\makeatletter
\makeatother
\makeatletter
\@ifpackageloaded{caption}{}{\usepackage{caption}}
\@ifpackageloaded{subcaption}{}{\usepackage{subcaption}}
\makeatother

\usepackage{bookmark}

\IfFileExists{xurl.sty}{\usepackage{xurl}}{} % add URL line breaks if available
\urlstyle{same} % disable monospaced font for URLs
\hypersetup{
  pdftitle={01\_Accession\_phenotypes},
  pdfauthor={Henry Janse van Rensburg},
  colorlinks=true,
  linkcolor={blue},
  filecolor={Maroon},
  citecolor={Blue},
  urlcolor={Blue},
  pdfcreator={LaTeX via pandoc}}


\title{01\_Accession\_phenotypes}
\author{Henry Janse van Rensburg}
\date{2025-02-05}

\begin{document}
\maketitle

\renewcommand*\contentsname{Table of contents}
{
\hypersetup{linkcolor=}
\setcounter{tocdepth}{2}
\tableofcontents
}

\section{Background}\label{background}

We conducted this experiment to study the phenotypic diversity within
\emph{Arabidopsis} accession to BX-growth feedbacks. We screened 410
accessions for their growth feedback on two different soils, BXplus and
BXminus. The phenotypes were measured twice a day for 20 days after
germination.

Materials and methodsWe grew 410 Arabidopsis accessions in triplicate on
BXplus and BXminus soils and phenotyped them twice a day. We extracted
the rosette area using ARADEEPOPSIS. The list of accessions used in this
study is available in the file \texttt{bx\_experiment\_genotypes.csv} in
the input folder.

\textbf{Purpose}: We performed this study to eventually perform GWAS on
the phenotypes.

\textbf{Important results:} We found that the phenotypes were
significantly different between the two soil types, and varied based on
plant genotype.

\textbf{Conclusions:} We concluded that the phenotypic diversity within
Arabidopsis accessions is significant and can be used for further
genetic studies.

\section{Setup}\label{setup}

\begin{Shaded}
\begin{Highlighting}[]
\FunctionTok{cat}\NormalTok{(}\StringTok{"}\SpecialCharTok{\textbackslash{}\textbackslash{}}\StringTok{small"}\NormalTok{)}
\end{Highlighting}
\end{Shaded}

\begin{verbatim}
\small
\end{verbatim}

\begin{Shaded}
\begin{Highlighting}[]
\NormalTok{knitr}\SpecialCharTok{::}\NormalTok{knit\_hooks}\SpecialCharTok{$}\FunctionTok{set}\NormalTok{(}\AttributeTok{purl =}\NormalTok{ knitr}\SpecialCharTok{::}\NormalTok{hook\_purl)}
\NormalTok{knitr}\SpecialCharTok{::}\NormalTok{opts\_chunk}\SpecialCharTok{$}\FunctionTok{set}\NormalTok{(}\AttributeTok{echo =} \ConstantTok{TRUE}\NormalTok{)}

\DocumentationTok{\#\#clear the objects from memory}
\FunctionTok{rm}\NormalTok{(}\AttributeTok{list=}\FunctionTok{ls}\NormalTok{())}
\CommentTok{\# Purpose: Clears all objects from the R environment to }
\CommentTok{\#ensure a clean workspace.}
\end{Highlighting}
\end{Shaded}

\subsection{Load libraries}\label{load-libraries}

\begin{Shaded}
\begin{Highlighting}[]
\CommentTok{\#| label: Libraries}
\CommentTok{\#| message: false}
\CommentTok{\#| warning: false}
\CommentTok{\#| paged{-}print: false}
\FunctionTok{library}\NormalTok{(tidyverse)}
\CommentTok{\#devtools::install\_github("Gregor{-}Mendel{-}Institute/gwaR"); }
\CommentTok{\#https://github.com/Gregor{-}Mendel{-}Institute/gwaR}
\FunctionTok{library}\NormalTok{(gwaR)}
\FunctionTok{library}\NormalTok{(magrittr)}
\FunctionTok{library}\NormalTok{(dplyr)}
\FunctionTok{library}\NormalTok{(ggplot2)}
\FunctionTok{library}\NormalTok{(tidyr)}
\FunctionTok{library}\NormalTok{(DescTools)}
\FunctionTok{library}\NormalTok{(ggpubr)}
\FunctionTok{library}\NormalTok{(MESS)}
\FunctionTok{library}\NormalTok{(cowplot)}
\FunctionTok{library}\NormalTok{(knitr)}
\FunctionTok{library}\NormalTok{(ggdendro)}
\ControlFlowTok{if}\NormalTok{ (}\SpecialCharTok{!}\FunctionTok{requireNamespace}\NormalTok{(}\StringTok{"BiocManager"}\NormalTok{, }\AttributeTok{quietly =} \ConstantTok{TRUE}\NormalTok{))}
    \FunctionTok{install.packages}\NormalTok{(}\StringTok{"BiocManager"}\NormalTok{)}
\NormalTok{BiocManager}\SpecialCharTok{::}\FunctionTok{install}\NormalTok{(}\StringTok{"ComplexHeatmap"}\NormalTok{)}
\FunctionTok{library}\NormalTok{(ComplexHeatmap)}
\FunctionTok{library}\NormalTok{(circlize)}
\FunctionTok{library}\NormalTok{(openxlsx)}
\CommentTok{\# Purpose: Loads the necessary R packages for data manipulation, visualization, }
\CommentTok{\#and GWAS analysis.}
\end{Highlighting}
\end{Shaded}

\subsection{Import phenotypic data}\label{import-phenotypic-data}

\textbf{Purpose}: Imports the phenotype data as
\texttt{plant\_region\_areas}and the SNPs identifid in the GWAs as
\texttt{tot\_plant\_specific\_snps}from CSV files into R as data frames.

\begin{Shaded}
\begin{Highlighting}[]
\FunctionTok{library}\NormalTok{(}\StringTok{"readxl"}\NormalTok{)}
\NormalTok{tot\_plant\_specific\_snps }\OtherTok{\textless{}{-}} \FunctionTok{read\_excel}\NormalTok{(}\StringTok{"Input/Dataset1.xlsx"}\NormalTok{, }\AttributeTok{sheet =} \StringTok{"SNP\_table"}\NormalTok{)}

\CommentTok{\# Change the orde of the table to show highest log10\_p values first}
\NormalTok{tot\_plant\_specific\_snps }\OtherTok{\textless{}{-}}\NormalTok{ tot\_plant\_specific\_snps[}\FunctionTok{order}\NormalTok{(tot\_plant\_specific\_snps}\SpecialCharTok{$}\NormalTok{log10\_p,}
    \AttributeTok{decreasing =} \ConstantTok{TRUE}\NormalTok{), ]}
\FunctionTok{kable}\NormalTok{(tot\_plant\_specific\_snps[}\DecValTok{1}\SpecialCharTok{:}\DecValTok{3}\NormalTok{, ], }\AttributeTok{caption =} \StringTok{"SNP table."}\NormalTok{)}
\end{Highlighting}
\end{Shaded}

\begin{longtable}[]{@{}rrrrlrr@{}}
\caption{SNP table.}\tabularnewline
\toprule\noalign{}
chrom & pos & log10\_p & pv & Significant & mac & dag \\
\midrule\noalign{}
\endfirsthead
\toprule\noalign{}
chrom & pos & log10\_p & pv & Significant & mac & dag \\
\midrule\noalign{}
\endhead
\bottomrule\noalign{}
\endlastfoot
3 & 21437549 & 8.990193 & 0 & Bonferroni & 5 & 13 \\
3 & 21437549 & 8.952353 & 0 & Bonferroni & 5 & 12 \\
3 & 21437549 & 8.870501 & 0 & Bonferroni & 5 & 14 \\
\end{longtable}

\begin{Shaded}
\begin{Highlighting}[]
\CommentTok{\# Import plant areas of accessions}
\NormalTok{plant\_region\_areas }\OtherTok{\textless{}{-}} \FunctionTok{read\_excel}\NormalTok{(}\StringTok{"Input/Dataset1.xlsx"}\NormalTok{, }\AttributeTok{sheet =} \StringTok{"plant\_region\_area"}\NormalTok{)}

\CommentTok{\# Purpose: Imports phenotype data files into R as data frames.}
\end{Highlighting}
\end{Shaded}

\section{Phenotypes}\label{phenotypes}

\subsection{Extract Phenotypes for SNP Rank
7}\label{extract-phenotypes-for-snp-rank-7}

\textbf{Purpose}: Extracts and combines phenotype data based on SNP rank
7 and prepares it as the \texttt{plant\_region\_areas\_SNP} dataframe in
for further analysis.

\begin{Shaded}
\begin{Highlighting}[]
\CommentTok{\# Extract plant area}
\NormalTok{temp1 }\OtherTok{\textless{}{-}}\NormalTok{ tot\_plant\_specific\_snps }\SpecialCharTok{\%\textgreater{}\%}
\NormalTok{  \{gwaR}\SpecialCharTok{::}\FunctionTok{phenotype\_by\_snp}\NormalTok{(plant\_region\_areas }\SpecialCharTok{\%\textgreater{}\%} \FunctionTok{mutate}\NormalTok{(}\AttributeTok{ACC\_ID =}\NormalTok{ genotype\_id),}
                          \StringTok{"plant\_region\_area"}\NormalTok{,}
                          \AttributeTok{gwas\_table =}\NormalTok{ .,}
                          \AttributeTok{specific =} \StringTok{"bx"}\NormalTok{,}
                          \AttributeTok{SNPrank =} \DecValTok{7}\NormalTok{,}
                          \AttributeTok{plot =} \ConstantTok{FALSE}\NormalTok{, }
                          \AttributeTok{SNPmatrix =} \StringTok{"Input/SNPmatrix\_full.fst"}\NormalTok{, }\CommentTok{\# Change this path!}
                          \AttributeTok{nobees =} \ConstantTok{TRUE}\NormalTok{)\}}

\CommentTok{\# Extract dag}
\NormalTok{temp2 }\OtherTok{\textless{}{-}}\NormalTok{ tot\_plant\_specific\_snps }\SpecialCharTok{\%\textgreater{}\%} 
\NormalTok{  \{gwaR}\SpecialCharTok{::}\FunctionTok{phenotype\_by\_snp}\NormalTok{(plant\_region\_areas }\SpecialCharTok{\%\textgreater{}\%} \FunctionTok{mutate}\NormalTok{(}\AttributeTok{ACC\_ID =}\NormalTok{ genotype\_id),}
                          \StringTok{"dag"}\NormalTok{,}
                          \AttributeTok{gwas\_table =}\NormalTok{ .,}
                          \AttributeTok{specific =} \StringTok{"bx"}\NormalTok{,}
                          \AttributeTok{SNPrank =} \DecValTok{7}\NormalTok{,}
                          \AttributeTok{plot =} \ConstantTok{FALSE}\NormalTok{, }
                          \AttributeTok{SNPmatrix =} \StringTok{"Input/SNPmatrix\_full.fst"}\NormalTok{, }\CommentTok{\# Change this path!}
                          \AttributeTok{nobees =} \ConstantTok{TRUE}\NormalTok{)\}}

\CommentTok{\# Combine to one df}
\NormalTok{plant\_region\_areas\_SNP }\OtherTok{\textless{}{-}} \FunctionTok{data.frame}\NormalTok{(}
  \AttributeTok{ACC\_ID =}\NormalTok{ temp1}\SpecialCharTok{$}\NormalTok{ACC\_ID, }
  \AttributeTok{bx =}\NormalTok{ temp1}\SpecialCharTok{$}\NormalTok{bx, }
  \AttributeTok{plant\_region\_area =}\NormalTok{ temp1}\SpecialCharTok{$}\NormalTok{phenotype\_value,}
  \AttributeTok{dag =} \FunctionTok{as.integer}\NormalTok{(temp2}\SpecialCharTok{$}\NormalTok{phenotype\_value),  }\CommentTok{\# Convert dag to integers}
  \AttributeTok{SNP =}\NormalTok{ temp1}\SpecialCharTok{$}\NormalTok{SNP}
\NormalTok{)}

\CommentTok{\# Remove temporary objects}
\FunctionTok{rm}\NormalTok{(temp1, temp2)}

\CommentTok{\# Purpose: Extracts and combines phenotype data based on SNP rank 7 and prepares it for further analysis.}
\FunctionTok{kable}\NormalTok{(plant\_region\_areas\_SNP[}\DecValTok{1}\SpecialCharTok{:}\DecValTok{3}\NormalTok{, ], }\AttributeTok{caption =} \StringTok{"Phenotype table for SNP rank 7."}\NormalTok{)}
\end{Highlighting}
\end{Shaded}

\begin{longtable}[]{@{}rlrrl@{}}
\caption{Phenotype table for SNP rank 7.}\tabularnewline
\toprule\noalign{}
ACC\_ID & bx & plant\_region\_area & dag & SNP \\
\midrule\noalign{}
\endfirsthead
\toprule\noalign{}
ACC\_ID & bx & plant\_region\_area & dag & SNP \\
\midrule\noalign{}
\endhead
\bottomrule\noalign{}
\endlastfoot
9830 & + & 636 & 1 & FALSE \\
6092 & + & 413 & 1 & FALSE \\
6022 & + & 380 & 1 & FALSE \\
\end{longtable}

\subsection{Extract phenotypes for DAG
20}\label{extract-phenotypes-for-dag-20}

\textbf{Purpose}: Create a new dataframe,
\texttt{plant\_region\_areas\_DAG20}, filters the data to include only
the first 20 days (DAG ≤ 20) for further analysis, calculates the mean
plant region area for each combination of factors, removes duplicates
and irrelevant columns, pivots the data to spread `bx' treatment values
into separate columns, calculates the area difference, and ranks the
genotypes by the area difference.

\begin{Shaded}
\begin{Highlighting}[]
\CommentTok{\# Filter the plant region areas for DAG ≤ 20}
\NormalTok{plant\_region\_areas\_DAG20 }\OtherTok{\textless{}{-}} \FunctionTok{subset}\NormalTok{(plant\_region\_areas, }\AttributeTok{subset =}\NormalTok{ dag }\SpecialCharTok{\textless{}=} \DecValTok{21}\NormalTok{)}

\CommentTok{\# Calculate mean plant region area for each combination of bx, dag, and genotype\_id}
\NormalTok{plant\_region\_areas\_DAG20 }\OtherTok{\textless{}{-}}\NormalTok{ plant\_region\_areas\_DAG20 }\SpecialCharTok{\%\textgreater{}\%}
  \FunctionTok{group\_by}\NormalTok{(bx, dag, genotype\_id) }\SpecialCharTok{\%\textgreater{}\%}
  \FunctionTok{mutate}\NormalTok{(}\AttributeTok{mean =} \FunctionTok{mean}\NormalTok{(plant\_region\_area, }\AttributeTok{na.rm =} \ConstantTok{TRUE}\NormalTok{)) }\SpecialCharTok{\%\textgreater{}\%}
  \FunctionTok{ungroup}\NormalTok{()}

\CommentTok{\# Remove duplicates and unnecessary columns}
\NormalTok{plant\_region\_areas\_DAG20 }\OtherTok{\textless{}{-}}\NormalTok{ plant\_region\_areas\_DAG20 }\SpecialCharTok{\%\textgreater{}\%}
  \FunctionTok{select}\NormalTok{(}\SpecialCharTok{{-}}\FunctionTok{c}\NormalTok{(}\DecValTok{1}\NormalTok{, }\DecValTok{6}\NormalTok{)) }\SpecialCharTok{\%\textgreater{}\%}
  \FunctionTok{distinct}\NormalTok{()}

\CommentTok{\# Pivot data to spread \textquotesingle{}bx\textquotesingle{} treatment values into separate columns and calculate the difference}
\NormalTok{plant\_region\_areas\_DAG20 }\OtherTok{\textless{}{-}}\NormalTok{ plant\_region\_areas\_DAG20 }\SpecialCharTok{\%\textgreater{}\%}
  \FunctionTok{spread}\NormalTok{(}\AttributeTok{key =} \StringTok{"bx"}\NormalTok{, }\AttributeTok{value =} \StringTok{"mean"}\NormalTok{) }\SpecialCharTok{\%\textgreater{}\%}
  \FunctionTok{mutate}\NormalTok{(}\AttributeTok{area\_diff =} \StringTok{\textasciigrave{}}\AttributeTok{+}\StringTok{\textasciigrave{}} \SpecialCharTok{{-}} \StringTok{\textasciigrave{}}\AttributeTok{{-}}\StringTok{\textasciigrave{}}\NormalTok{) }\SpecialCharTok{\%\textgreater{}\%}
  \FunctionTok{arrange}\NormalTok{(}\FunctionTok{desc}\NormalTok{(area\_diff))}

\CommentTok{\# Rank the entries by area difference}
\NormalTok{plant\_region\_areas\_DAG20 }\OtherTok{\textless{}{-}}\NormalTok{ plant\_region\_areas\_DAG20 }\SpecialCharTok{\%\textgreater{}\%}
  \FunctionTok{mutate}\NormalTok{(}\AttributeTok{rank =} \FunctionTok{row\_number}\NormalTok{(}\FunctionTok{desc}\NormalTok{(area\_diff)))}

\CommentTok{\# Purpose: Ranks the genotypes by the calculated area difference.}
\FunctionTok{kable}\NormalTok{(plant\_region\_areas\_DAG20[}\DecValTok{1}\SpecialCharTok{:}\DecValTok{3}\NormalTok{, ], }\AttributeTok{caption =} \StringTok{"Phenotype table for DAG 20."}\NormalTok{)}
\end{Highlighting}
\end{Shaded}

\begin{longtable}[]{@{}lrrrrrr@{}}
\caption{Phenotype table for DAG 20.}\tabularnewline
\toprule\noalign{}
genotype & genotype\_id & dag & - & + & area\_diff & rank \\
\midrule\noalign{}
\endfirsthead
\toprule\noalign{}
genotype & genotype\_id & dag & - & + & area\_diff & rank \\
\midrule\noalign{}
\endhead
\bottomrule\noalign{}
\endlastfoot
DraIV 3-7 & 5921 & 21 & 29516.83 & 115727.00 & 86210.17 & 1 \\
DraIV 3-7 & 5921 & 20 & 23975.00 & 98518.33 & 74543.33 & 2 \\
Hov1-10 & 6035 & 21 & 15445.25 & 86467.33 & 71022.08 & 3 \\
\end{longtable}

\subsection{Assign BX feedback and Prepare Data for
Heatmap}\label{assign-bx-feedback-and-prepare-data-for-heatmap}

\textbf{Purpose}: Assigns a preference (``plus'' or ``minus'') to each
genotype based on the area difference at day 10 and merges the
preference data back into the main dataframe. We chose DAG10, as this is
where SNPMMF1 showed up in the GWAS analysis.

\begin{Shaded}
\begin{Highlighting}[]
\CommentTok{\# Step 1: Create a dataframe with preferences based on "dag" 15}
\NormalTok{preferences }\OtherTok{\textless{}{-}}\NormalTok{ plant\_region\_areas\_DAG20 }\SpecialCharTok{\%\textgreater{}\%}
  \FunctionTok{filter}\NormalTok{(dag }\SpecialCharTok{==} \DecValTok{10}\NormalTok{) }\SpecialCharTok{\%\textgreater{}\%}
  \FunctionTok{mutate}\NormalTok{(}\AttributeTok{preference =} \FunctionTok{case\_when}\NormalTok{(}
\NormalTok{    area\_diff }\SpecialCharTok{\textgreater{}=} \DecValTok{0} \SpecialCharTok{\textasciitilde{}} \StringTok{"plus"}\NormalTok{,}
    \ConstantTok{TRUE} \SpecialCharTok{\textasciitilde{}} \StringTok{"minus"}
\NormalTok{  )) }\SpecialCharTok{\%\textgreater{}\%}
  \FunctionTok{select}\NormalTok{(genotype, preference)}
\CommentTok{\# Purpose: Assigns a preference ("plus" or "minus") to each genotype based on the area difference at day 10.}

\CommentTok{\# Step 2: Merge preferences back to the original dataframe}
\NormalTok{plant\_region\_areas\_DAG20 }\OtherTok{\textless{}{-}}\NormalTok{ plant\_region\_areas\_DAG20 }\SpecialCharTok{\%\textgreater{}\%}
  \FunctionTok{left\_join}\NormalTok{(preferences, }\AttributeTok{by =} \StringTok{"genotype"}\NormalTok{)}
\CommentTok{\# Purpose: Merges the preference data back into the main dataframe.}

\CommentTok{\# Order by rank (if needed, depending on the requirement)}
\NormalTok{plant\_region\_areas\_DAG20 }\OtherTok{\textless{}{-}}\NormalTok{ plant\_region\_areas\_DAG20 }\SpecialCharTok{\%\textgreater{}\%}
  \FunctionTok{arrange}\NormalTok{(}\FunctionTok{desc}\NormalTok{(rank))}
\CommentTok{\# Purpose: Orders the dataframe by rank for consistent visualization.}
\end{Highlighting}
\end{Shaded}

\subsection{Prepare Data for Heatmap for All
Accessions}\label{prepare-data-for-heatmap-for-all-accessions}

\textbf{Purpose}: Prepares the genotype and time order for the heatmap.

\begin{Shaded}
\begin{Highlighting}[]
\NormalTok{genotype\_DAG20 }\OtherTok{\textless{}{-}}\NormalTok{ plant\_region\_areas\_DAG20 }\SpecialCharTok{\%\textgreater{}\%} \FunctionTok{filter}\NormalTok{(dag }\SpecialCharTok{==} \DecValTok{20}\NormalTok{)}\CommentTok{\# dag 20 is the 20th day after germination was picked up}
\NormalTok{genotype\_order }\OtherTok{\textless{}{-}}\NormalTok{ genotype\_DAG20}\SpecialCharTok{$}\NormalTok{genotype}
\NormalTok{time\_order }\OtherTok{\textless{}{-}} \DecValTok{0}\SpecialCharTok{:}\DecValTok{20}
\end{Highlighting}
\end{Shaded}

\subsection{Plot Growth Over Time for All
Accessions}\label{plot-growth-over-time-for-all-accessions}

\textbf{Purpose}: Creates a heatmap showing the area differences over
time for all accessions, grouped by preference.

\subsection{Top 10 Positive and Negative
Accessions}\label{top-10-positive-and-negative-accessions}

\textbf{Purpose}: Filters the data from (\textbf{?@fig-heatmap-ACCs})
include only the top 20 positive and negative accessions based on area
difference and applies a pseudo-logarithmic transformation to the area
difference for better visualization at early timepoints.

\begin{Shaded}
\begin{Highlighting}[]
\CommentTok{\# Filter for specific genotypes with the highest opposing feedbacks}
\CommentTok{\#and days less than or equal to 20}

\CommentTok{\# Filter top 20 positive and negative accessions based on area difference}
\NormalTok{plant\_region\_areas\_DAG20\_extreme20 }\OtherTok{\textless{}{-}}\NormalTok{ plant\_region\_areas\_DAG20 }\SpecialCharTok{\%\textgreater{}\%}
  \FunctionTok{filter}\NormalTok{(genotype\_id }\SpecialCharTok{\%in\%} \FunctionTok{c}\NormalTok{(}\StringTok{"9851"}\NormalTok{, }\StringTok{"7350"}\NormalTok{, }\StringTok{"5867"}\NormalTok{, }\StringTok{"5921"}\NormalTok{, }\StringTok{"9972"}\NormalTok{, }\StringTok{"10006"}\NormalTok{, }
                            \StringTok{"9823"}\NormalTok{, }\StringTok{"7161"}\NormalTok{, }\StringTok{"6088"}\NormalTok{, }\StringTok{"9084"}\NormalTok{, }\StringTok{"6909"}\NormalTok{, }\StringTok{"9392"}\NormalTok{, }
                            \StringTok{"8365"}\NormalTok{, }\StringTok{"6136"}\NormalTok{, }\StringTok{"7217"}\NormalTok{, }\StringTok{"506"}\NormalTok{, }\StringTok{"7520"}\NormalTok{, }\StringTok{"6108"}\NormalTok{, }
                            \StringTok{"9653"}\NormalTok{, }\StringTok{"9821"}\NormalTok{, }\StringTok{"9591"}\NormalTok{),}
\NormalTok{         dag }\SpecialCharTok{\textless{}=} \DecValTok{20}\NormalTok{) }\SpecialCharTok{\%\textgreater{}\%}
  \FunctionTok{mutate}\NormalTok{(}\AttributeTok{group =} \FunctionTok{factor}\NormalTok{(genotype, }
                        \AttributeTok{levels =} \FunctionTok{c}\NormalTok{(}\StringTok{\textquotesingle{}IP{-}Hue{-}3\textquotesingle{}}\NormalTok{, }\StringTok{\textquotesingle{}Tac{-}0\textquotesingle{}}\NormalTok{, }\StringTok{\textquotesingle{}DraIV 3{-}7\textquotesingle{}}\NormalTok{, }\StringTok{\textquotesingle{}Dra2{-}1\textquotesingle{}}\NormalTok{, }
                                   \StringTok{\textquotesingle{}Bozen{-}1.2\textquotesingle{}}\NormalTok{, }\StringTok{\textquotesingle{}Kastel{-}1\textquotesingle{}}\NormalTok{, }\StringTok{\textquotesingle{}IP{-}Bae{-}0\textquotesingle{}}\NormalTok{, }\StringTok{\textquotesingle{}Gd{-}1\textquotesingle{}}\NormalTok{, }
                                   \StringTok{\textquotesingle{}Stu1{-}1\textquotesingle{}}\NormalTok{, }\StringTok{\textquotesingle{}Lerik2{-}6\textquotesingle{}}\NormalTok{, }\StringTok{"Col{-}0"}\NormalTok{, }\StringTok{"Hadd{-}3"}\NormalTok{, }
                                   \StringTok{"Rak{-}2"}\NormalTok{, }\StringTok{"T840"}\NormalTok{, }\StringTok{"Lm{-}2"}\NormalTok{, }\StringTok{"BRR60"}\NormalTok{, }\StringTok{"Lp2{-}2"}\NormalTok{, }
                                   \StringTok{"T480"}\NormalTok{, }\StringTok{"Giffo{-}1"}\NormalTok{, }\StringTok{"IP{-}Aru{-}0"}\NormalTok{, }\StringTok{"IP{-}Vad{-}0"}\NormalTok{)))}

\CommentTok{\# Pseudo{-}log10 transformation function}
\NormalTok{pseudoLog10 }\OtherTok{\textless{}{-}} \ControlFlowTok{function}\NormalTok{(x) \{ }
  \FunctionTok{log}\NormalTok{((x }\SpecialCharTok{/} \DecValTok{2}\NormalTok{) }\SpecialCharTok{+} \FunctionTok{sqrt}\NormalTok{((x }\SpecialCharTok{/} \DecValTok{2}\NormalTok{)}\SpecialCharTok{\^{}}\DecValTok{2} \SpecialCharTok{+} \DecValTok{1}\NormalTok{)) }\SpecialCharTok{/} \FunctionTok{log}\NormalTok{(}\DecValTok{10}\NormalTok{)}
\NormalTok{\}}

\CommentTok{\# Apply transformation}
\NormalTok{plant\_region\_areas\_DAG20\_extreme20 }\OtherTok{\textless{}{-}}\NormalTok{ plant\_region\_areas\_DAG20\_extreme20 }\SpecialCharTok{\%\textgreater{}\%}
  \FunctionTok{mutate}\NormalTok{(}\AttributeTok{pseudo\_log\_area\_diff =} \FunctionTok{pseudoLog10}\NormalTok{(area\_diff))}
\end{Highlighting}
\end{Shaded}

\subsection{Hierarchical Clustering and Heatmap
Preparation}\label{hierarchical-clustering-and-heatmap-preparation}

\textbf{Purpos}e: Transforms the data into a wide format suitable for
clustering and performs hierarchical clustering on the distance matrix.
The reordered groups are then used to arrange the data frame for
visualization. Creates and displays a heatmap showing the growth
feedback over time for the top 20 extreme accessions.

\subsection{Hierarchical Clustering: Dynamic Time Warping
(DTW)}\label{hierarchical-clustering-dynamic-time-warping-dtw}

\begin{Shaded}
\begin{Highlighting}[]
\CommentTok{\# Load required packages}
\ControlFlowTok{if}\NormalTok{ (}\SpecialCharTok{!}\FunctionTok{requireNamespace}\NormalTok{(}\StringTok{"dtw"}\NormalTok{, }\AttributeTok{quietly =} \ConstantTok{TRUE}\NormalTok{)) }\FunctionTok{install.packages}\NormalTok{(}\StringTok{"dtw"}\NormalTok{)}
\ControlFlowTok{if}\NormalTok{ (}\SpecialCharTok{!}\FunctionTok{requireNamespace}\NormalTok{(}\StringTok{"ComplexHeatmap"}\NormalTok{, }\AttributeTok{quietly =} \ConstantTok{TRUE}\NormalTok{)) }\FunctionTok{install.packages}\NormalTok{(}\StringTok{"ComplexHeatmap"}\NormalTok{)}
\FunctionTok{library}\NormalTok{(dtw)}
\FunctionTok{library}\NormalTok{(ComplexHeatmap)}
\FunctionTok{library}\NormalTok{(circlize)}
\FunctionTok{library}\NormalTok{(dplyr)}
\FunctionTok{library}\NormalTok{(tidyr)}
\end{Highlighting}
\end{Shaded}

\subsubsection{Prepare Data}\label{prepare-data}

\begin{Shaded}
\begin{Highlighting}[]
\CommentTok{\# Transform data into wide format}
\NormalTok{wide\_data }\OtherTok{\textless{}{-}}\NormalTok{ plant\_region\_areas\_DAG20\_extreme20 }\SpecialCharTok{\%\textgreater{}\%}
  \FunctionTok{select}\NormalTok{(dag, group, pseudo\_log\_area\_diff) }\SpecialCharTok{\%\textgreater{}\%}
  \FunctionTok{pivot\_wider}\NormalTok{(}\AttributeTok{names\_from =}\NormalTok{ dag, }\AttributeTok{values\_from =}\NormalTok{ pseudo\_log\_area\_diff)}

\CommentTok{\# Ensure columns (days after germination) are sorted numerically}
\NormalTok{col\_order }\OtherTok{\textless{}{-}} \FunctionTok{order}\NormalTok{(}\FunctionTok{as.numeric}\NormalTok{(}\FunctionTok{colnames}\NormalTok{(wide\_data)[}\SpecialCharTok{{-}}\DecValTok{1}\NormalTok{]))  }\CommentTok{\# Exclude \textquotesingle{}group\textquotesingle{} column}
\NormalTok{wide\_data }\OtherTok{\textless{}{-}}\NormalTok{ wide\_data }\SpecialCharTok{\%\textgreater{}\%}
  \FunctionTok{select}\NormalTok{(group, }\FunctionTok{colnames}\NormalTok{(wide\_data)[}\SpecialCharTok{{-}}\DecValTok{1}\NormalTok{][col\_order])}

\CommentTok{\# Convert to matrix}
\NormalTok{heatmap\_matrix }\OtherTok{\textless{}{-}} \FunctionTok{as.matrix}\NormalTok{(wide\_data[, }\SpecialCharTok{{-}}\DecValTok{1}\NormalTok{])  }\CommentTok{\# Exclude \textquotesingle{}group\textquotesingle{}}
\FunctionTok{rownames}\NormalTok{(heatmap\_matrix) }\OtherTok{\textless{}{-}}\NormalTok{ wide\_data}\SpecialCharTok{$}\NormalTok{group   }\CommentTok{\# Set row names to genotypes}

\CommentTok{\# **2. Calculate DTW Distance Matrix**}
\CommentTok{\# Using vectorized computation}
\NormalTok{dtw\_dist }\OtherTok{\textless{}{-}}\NormalTok{ proxy}\SpecialCharTok{::}\FunctionTok{dist}\NormalTok{(heatmap\_matrix, }\AttributeTok{method =} \ControlFlowTok{function}\NormalTok{(x, y) }\FunctionTok{dtw}\NormalTok{(x, y)}\SpecialCharTok{$}\NormalTok{distance)}
\end{Highlighting}
\end{Shaded}

\subsubsection{Perform Hierarchical
Clustering}\label{perform-hierarchical-clustering}

\begin{Shaded}
\begin{Highlighting}[]
\NormalTok{fit }\OtherTok{\textless{}{-}} \FunctionTok{hclust}\NormalTok{(dtw\_dist, }\AttributeTok{method =} \StringTok{"ward.D2"}\NormalTok{)}

\CommentTok{\# Reverse row order for visualization (optional)}
\NormalTok{row\_order }\OtherTok{\textless{}{-}} \FunctionTok{rev}\NormalTok{(}\DecValTok{1}\SpecialCharTok{:}\FunctionTok{nrow}\NormalTok{(heatmap\_matrix))}
\end{Highlighting}
\end{Shaded}

\subsubsection{Heatmap}\label{heatmap}

\begin{Shaded}
\begin{Highlighting}[]
\CommentTok{\# **4. Create and Customize the Heatmap**}
\NormalTok{ht }\OtherTok{\textless{}{-}} \FunctionTok{Heatmap}\NormalTok{(}
\NormalTok{  heatmap\_matrix,}
  \AttributeTok{name =} \StringTok{"Phenotype Change"}\NormalTok{,}
  \AttributeTok{col =} \FunctionTok{colorRamp2}\NormalTok{(}
    \FunctionTok{c}\NormalTok{(}\FunctionTok{min}\NormalTok{(heatmap\_matrix), }\DecValTok{0}\NormalTok{, }\FunctionTok{max}\NormalTok{(heatmap\_matrix)),}
    \FunctionTok{c}\NormalTok{(}\StringTok{"orange2"}\NormalTok{, }\StringTok{"white"}\NormalTok{, }\StringTok{"\#0072B2"}\NormalTok{),}
\NormalTok{  ),}
  \AttributeTok{cluster\_rows =}\NormalTok{ fit,}
  \AttributeTok{cluster\_columns =} \ConstantTok{FALSE}\NormalTok{,}
  \AttributeTok{row\_order =}\NormalTok{ row\_order,}
  \AttributeTok{show\_row\_dend =} \ConstantTok{TRUE}\NormalTok{,}
  \AttributeTok{row\_dend\_side =} \StringTok{"right"}\NormalTok{,}
  \AttributeTok{row\_names\_side =} \StringTok{"left"}\NormalTok{,}
  \AttributeTok{row\_title =} \StringTok{""}\NormalTok{,}
  \AttributeTok{row\_dend\_width =} \FunctionTok{unit}\NormalTok{(}\DecValTok{30}\NormalTok{, }\StringTok{"mm"}\NormalTok{),}
  \AttributeTok{row\_names\_gp =} \FunctionTok{gpar}\NormalTok{(}\AttributeTok{fontsize =} \DecValTok{7}\NormalTok{),}
  \AttributeTok{column\_names\_gp =} \FunctionTok{gpar}\NormalTok{(}\AttributeTok{fontsize =} \DecValTok{7}\NormalTok{),}
  \AttributeTok{column\_title =} \StringTok{"Days after germination"}\NormalTok{,}
  \AttributeTok{column\_title\_gp =} \FunctionTok{gpar}\NormalTok{(}\AttributeTok{fontsize =} \DecValTok{7}\NormalTok{),}
  \AttributeTok{column\_title\_side =} \StringTok{"bottom"}\NormalTok{,}
  \AttributeTok{column\_names\_rot =} \DecValTok{0}\NormalTok{,}
  \AttributeTok{heatmap\_legend\_param =} \FunctionTok{list}\NormalTok{(}\AttributeTok{title =} \StringTok{"Log Area Diff"}\NormalTok{, }\AttributeTok{legend\_direction =} \StringTok{"horizontal"}\NormalTok{),}
  \AttributeTok{show\_heatmap\_legend =} \ConstantTok{TRUE}
\NormalTok{)}

\CommentTok{\# Draw the heatmap in the document}
\FunctionTok{draw}\NormalTok{(ht, }\AttributeTok{heatmap\_legend\_side =} \StringTok{"bottom"}\NormalTok{)}
\end{Highlighting}
\end{Shaded}

\begin{figure}[H]

{\centering \pandocbounded{\includegraphics[keepaspectratio]{01_Accession_phenotypes_files/figure-pdf/FigS1b-heatmap-ACCs20-1.pdf}}

}

\caption{Heatmap showing the evolution of growth feedbacks as Pseudo-log
transformed values over time for the top 20 accessions, grouped by
preference. Hierarchical clustering was performed using the Dynamic Time
Warping method to account for changes over time.}

\end{figure}%

\begin{Shaded}
\begin{Highlighting}[]
\CommentTok{\# **5. Export Heatmap to PDF**}
\NormalTok{output\_dir }\OtherTok{\textless{}{-}} \StringTok{"Output/Figures"}
\ControlFlowTok{if}\NormalTok{ (}\SpecialCharTok{!}\FunctionTok{dir.exists}\NormalTok{(output\_dir)) }\FunctionTok{dir.create}\NormalTok{(output\_dir, }\AttributeTok{recursive =} \ConstantTok{TRUE}\NormalTok{)}
\FunctionTok{pdf}\NormalTok{(}\FunctionTok{file.path}\NormalTok{(output\_dir, }\StringTok{"ACC20\_rel\_growth\_time\_DTW.pdf"}\NormalTok{), }\AttributeTok{width =} \DecValTok{6}\NormalTok{, }\AttributeTok{height =} \DecValTok{5}\NormalTok{)}
\FunctionTok{draw}\NormalTok{(ht, }\AttributeTok{heatmap\_legend\_side =} \StringTok{"bottom"}\NormalTok{)}
\FunctionTok{dev.off}\NormalTok{()}
\end{Highlighting}
\end{Shaded}

\begin{verbatim}
pdf 
  2 
\end{verbatim}

\subsection{Accession feedback
distribution}\label{accession-feedback-distribution}

\subsubsection{Prepare data}\label{prepare-data-1}

We filter the data for DAG 15, where SNPMMF1 was still showing up.

\textbf{Purpose}: Filters the data for day 15, assigns a preference
(``plus'' or ``minus'') based on area difference, and ranks the
accessions accordingly. We used dag 15 as this was the last day where
SNPMMF1 showed up in the GWAS analysis and where all the accessions were
still prsent.

\begin{Shaded}
\begin{Highlighting}[]
\CommentTok{\# Filter for observations on day 15, assign preference, and sort}
\NormalTok{plant\_region\_areas\_DAG15 }\OtherTok{\textless{}{-}}\NormalTok{ plant\_region\_areas\_DAG20 }\SpecialCharTok{\%\textgreater{}\%}
  \FunctionTok{filter}\NormalTok{(dag }\SpecialCharTok{==} \DecValTok{15}\NormalTok{) }\SpecialCharTok{\%\textgreater{}\%}
  \FunctionTok{mutate}\NormalTok{(}\AttributeTok{preference =} \FunctionTok{case\_when}\NormalTok{(}
\NormalTok{    area\_diff }\SpecialCharTok{\textgreater{}=} \DecValTok{0} \SpecialCharTok{\textasciitilde{}} \StringTok{"plus"}\NormalTok{,}
    \ConstantTok{TRUE} \SpecialCharTok{\textasciitilde{}} \StringTok{"minus"}
\NormalTok{  )) }\SpecialCharTok{\%\textgreater{}\%}
  \FunctionTok{arrange}\NormalTok{(}\FunctionTok{desc}\NormalTok{(area\_diff)) }\SpecialCharTok{\%\textgreater{}\%}
  \FunctionTok{mutate}\NormalTok{(}\AttributeTok{rank =} \FunctionTok{row\_number}\NormalTok{())}
\end{Highlighting}
\end{Shaded}

Plot relative growth feedbacks

\textbf{Purpose}: Creates a bar plot showing the relative growth
feedback for each accession at day 15, ranked and colored by preference.

\begin{Shaded}
\begin{Highlighting}[]
\CommentTok{\# Prepare data for specific genotypes to label}
\NormalTok{labels\_data }\OtherTok{\textless{}{-}}\NormalTok{ plant\_region\_areas\_DAG15 }\SpecialCharTok{\%\textgreater{}\%}
  \FunctionTok{filter}\NormalTok{(genotype }\SpecialCharTok{\%in\%} \FunctionTok{c}\NormalTok{(}\StringTok{"IP{-}Hue{-}3"}\NormalTok{, }\StringTok{"Col{-}0"}\NormalTok{, }\StringTok{"Hadd{-}3"}\NormalTok{)) }\SpecialCharTok{\%\textgreater{}\%}
  \FunctionTok{mutate}\NormalTok{(}
    \AttributeTok{text\_position =} \FunctionTok{ifelse}\NormalTok{(area\_diff }\SpecialCharTok{\textgreater{}} \DecValTok{0}\NormalTok{, area\_diff }\SpecialCharTok{+} \DecValTok{4000}\NormalTok{, area\_diff }\SpecialCharTok{{-}} \DecValTok{5000}\NormalTok{) }\CommentTok{\# Offset for labels}
\NormalTok{  )}

\CommentTok{\# Create the plot}
\NormalTok{ACC\_ratio }\OtherTok{\textless{}{-}} \FunctionTok{ggplot}\NormalTok{(plant\_region\_areas\_DAG15, }\FunctionTok{aes}\NormalTok{(}\AttributeTok{x =}\NormalTok{ rank, }\AttributeTok{y =}\NormalTok{ area\_diff, }\AttributeTok{fill =}\NormalTok{ preference)) }\SpecialCharTok{+}
  
  \CommentTok{\# Add bars with original fill color based on "preference"}
  \FunctionTok{geom\_bar}\NormalTok{(}\AttributeTok{stat =} \StringTok{"identity"}\NormalTok{, }\AttributeTok{linewidth =} \FloatTok{0.001}\NormalTok{, }\AttributeTok{alpha =} \FloatTok{0.8}\NormalTok{, }\AttributeTok{width =} \DecValTok{1}\NormalTok{) }\SpecialCharTok{+}
  
  \CommentTok{\# Add a second layer to highlight specific genotypes in black}
  \FunctionTok{geom\_bar}\NormalTok{(}
    \AttributeTok{data =}\NormalTok{ plant\_region\_areas\_DAG15 }\SpecialCharTok{\%\textgreater{}\%} \FunctionTok{filter}\NormalTok{(genotype }\SpecialCharTok{\%in\%} \FunctionTok{c}\NormalTok{(}\StringTok{"IP{-}Hue{-}3"}\NormalTok{, }\StringTok{"Col{-}0"}\NormalTok{, }\StringTok{"Hadd{-}3"}\NormalTok{)),}
    \FunctionTok{aes}\NormalTok{(}\AttributeTok{x =}\NormalTok{ rank, }\AttributeTok{y =}\NormalTok{ area\_diff),}
    \AttributeTok{stat =} \StringTok{"identity"}\NormalTok{, }\AttributeTok{fill =} \StringTok{"black"}\NormalTok{, }\AttributeTok{linewidth =} \FloatTok{0.001}\NormalTok{, }\AttributeTok{width =} \DecValTok{1}
\NormalTok{  ) }\SpecialCharTok{+}
  
  \CommentTok{\# Customize the fill scale for the legend}
  \FunctionTok{scale\_fill\_manual}\NormalTok{(}
    \AttributeTok{values =} \FunctionTok{c}\NormalTok{(}\StringTok{"orange2"}\NormalTok{, }\StringTok{"\#0072B2"}\NormalTok{),}
    \AttributeTok{name =} \StringTok{"BX feedback"}\NormalTok{,  }\CommentTok{\# Update legend title}
    \AttributeTok{labels =} \FunctionTok{c}\NormalTok{(}\StringTok{"negative"}\NormalTok{, }\StringTok{"positive"}\NormalTok{)  }\CommentTok{\# Update legend labels}
\NormalTok{  ) }\SpecialCharTok{+}
  
  \CommentTok{\# Add points at the exact value of area\_diff}
  \FunctionTok{geom\_point}\NormalTok{(}\AttributeTok{data =}\NormalTok{ labels\_data,}
             \FunctionTok{aes}\NormalTok{(}\AttributeTok{y =}\NormalTok{ area\_diff), }
             \AttributeTok{color =} \StringTok{"black"}\NormalTok{, }\AttributeTok{size =} \FloatTok{1.2}\NormalTok{) }\SpecialCharTok{+}
  
  \CommentTok{\# Add text labels dynamically positioned}
  \FunctionTok{geom\_text}\NormalTok{(}\AttributeTok{data =}\NormalTok{ labels\_data,}
            \FunctionTok{aes}\NormalTok{(}\AttributeTok{y =}\NormalTok{ text\_position, }\AttributeTok{label =}\NormalTok{ genotype), }
            \AttributeTok{color =} \StringTok{"black"}\NormalTok{, }\AttributeTok{size =} \DecValTok{3}\NormalTok{, }\AttributeTok{hjust =} \FloatTok{0.5}\NormalTok{, }\AttributeTok{vjust =} \SpecialCharTok{{-}}\FloatTok{0.2}\NormalTok{) }\SpecialCharTok{+}  \CommentTok{\# Increased text size}
  
  \FunctionTok{xlim}\NormalTok{(}\SpecialCharTok{{-}}\DecValTok{2}\NormalTok{, }\DecValTok{412}\NormalTok{) }\SpecialCharTok{+}
  
  \CommentTok{\# Modify the y{-}axis label and scale}
  \FunctionTok{labs}\NormalTok{(}
    \AttributeTok{y =} \FunctionTok{expression}\NormalTok{(}\StringTok{"Growth feedback (10"}\SpecialCharTok{\^{}}\DecValTok{4}\SpecialCharTok{*}\StringTok{")"}\NormalTok{), }
    \AttributeTok{x =} \StringTok{"Accession rank"}
\NormalTok{  ) }\SpecialCharTok{+}
  
  \CommentTok{\# Scale y{-}axis to divide values by 10,000 for display}
  \FunctionTok{scale\_y\_continuous}\NormalTok{(}\AttributeTok{labels =} \ControlFlowTok{function}\NormalTok{(x) x }\SpecialCharTok{/} \DecValTok{10000}\NormalTok{) }\SpecialCharTok{+}
  
  \FunctionTok{ggtitle}\NormalTok{(}\StringTok{"Growth feedback distribution of accessions"}\NormalTok{) }\SpecialCharTok{+}
  \FunctionTok{theme\_bw}\NormalTok{() }\SpecialCharTok{+}
  \FunctionTok{theme}\NormalTok{(}
    \AttributeTok{axis.text =} \FunctionTok{element\_text}\NormalTok{(}\AttributeTok{size =} \DecValTok{8}\NormalTok{),  }\CommentTok{\# Slightly increased font size}
    \AttributeTok{axis.title =} \FunctionTok{element\_text}\NormalTok{(}\AttributeTok{size =} \DecValTok{9}\NormalTok{), }
    \AttributeTok{plot.title =} \FunctionTok{element\_text}\NormalTok{(}\AttributeTok{size =} \DecValTok{9}\NormalTok{, }\AttributeTok{face =} \StringTok{"bold"}\NormalTok{),}
    \AttributeTok{legend.position =} \FunctionTok{c}\NormalTok{(}\FloatTok{0.95}\NormalTok{, }\FloatTok{0.95}\NormalTok{),  }\CommentTok{\# Custom position for the top{-}right corner}
    \AttributeTok{legend.justification =} \FunctionTok{c}\NormalTok{(}\DecValTok{1}\NormalTok{, }\DecValTok{1}\NormalTok{),   }\CommentTok{\# Align legend box to top{-}right}
    \AttributeTok{legend.background =} \FunctionTok{element\_rect}\NormalTok{(}\AttributeTok{fill =} \StringTok{"white"}\NormalTok{, }\AttributeTok{color =} \StringTok{"black"}\NormalTok{),  }\CommentTok{\# Optional: Add background}
    \AttributeTok{legend.title =} \FunctionTok{element\_text}\NormalTok{(}\AttributeTok{size =} \DecValTok{7}\NormalTok{),  }\CommentTok{\# Increased legend title size}
    \AttributeTok{legend.text =} \FunctionTok{element\_text}\NormalTok{(}\AttributeTok{size =} \DecValTok{6}\NormalTok{)    }\CommentTok{\# Increased legend items size}
\NormalTok{  )}

\CommentTok{\# Display the plot}
\NormalTok{ACC\_ratio}
\end{Highlighting}
\end{Shaded}

\begin{figure}[H]

{\centering \pandocbounded{\includegraphics[keepaspectratio]{01_Accession_phenotypes_files/figure-pdf/Fig1c-ACC-Distribution-1.pdf}}

}

\caption{Distribution of growth feedbackas for all 410 accesions used in
the GWAS. The most extreme accessions and Col-0 is indicated on the
plot.}

\end{figure}%

\begin{Shaded}
\begin{Highlighting}[]
\CommentTok{\# Save the plot}
\FunctionTok{ggsave}\NormalTok{(}\StringTok{"ACC\_distr\_rel\_feedback.pdf"}\NormalTok{, }\AttributeTok{plot =}\NormalTok{ ACC\_ratio, }\AttributeTok{width =} \DecValTok{16}\NormalTok{, }\AttributeTok{height =} \DecValTok{8}\NormalTok{, }\AttributeTok{units =} \StringTok{"cm"}\NormalTok{, }\AttributeTok{path =} \StringTok{"Output/Figures/"}\NormalTok{)}
\end{Highlighting}
\end{Shaded}

\subsection{Growth of extreme
accessions}\label{growth-of-extreme-accessions}

\subsubsection{IP-Hue-3}\label{ip-hue-3}

Filter for IP-Hue-3 and Plot the Growth Over Time

\textbf{Purpose}: Filters the data to include only the genotype with
genotype\_id 9851 (IP-Hue-3), calculates the mean plant region area for
each combination of bx and dag, removes duplicates and irrelevant
columns, pivots the data to spread `bx' treatment values into separate
columns, calculates the area difference, and ranks the genotypes by the
area difference.

\begin{Shaded}
\begin{Highlighting}[]
\CommentTok{\# Filter data for genotype\_id 9851 (IP{-}Hue{-}3) and DAG ≤ 20}
\NormalTok{plant\_region\_areas\_IPHue3 }\OtherTok{\textless{}{-}}\NormalTok{ plant\_region\_areas }\SpecialCharTok{\%\textgreater{}\%}
  \FunctionTok{filter}\NormalTok{(dag }\SpecialCharTok{\textless{}=} \DecValTok{20}\NormalTok{, genotype\_id }\SpecialCharTok{==} \DecValTok{9851}\NormalTok{)}

\CommentTok{\# Calculate mean plant region area for each bx and dag combination}
\NormalTok{plant\_region\_areas\_IPHue3 }\OtherTok{\textless{}{-}}\NormalTok{ plant\_region\_areas\_IPHue3 }\SpecialCharTok{\%\textgreater{}\%}
  \FunctionTok{group\_by}\NormalTok{(bx, dag) }\SpecialCharTok{\%\textgreater{}\%}
  \FunctionTok{summarise}\NormalTok{(}\AttributeTok{mean =} \FunctionTok{mean}\NormalTok{(plant\_region\_area, }\AttributeTok{na.rm =} \ConstantTok{TRUE}\NormalTok{), }\AttributeTok{.groups =} \StringTok{"drop"}\NormalTok{)}

\CommentTok{\# Remove unnecessary columns and ensure unique rows}
\NormalTok{plant\_region\_areas\_IPHue3 }\OtherTok{\textless{}{-}} \FunctionTok{distinct}\NormalTok{(plant\_region\_areas\_IPHue3)}

\CommentTok{\# Transform data: separate columns for each treatment}
\NormalTok{shaded\_data\_IP\_Hue }\OtherTok{\textless{}{-}}\NormalTok{ plant\_region\_areas\_IPHue3 }\SpecialCharTok{\%\textgreater{}\%}
  \FunctionTok{pivot\_wider}\NormalTok{(}\AttributeTok{names\_from =}\NormalTok{ bx, }\AttributeTok{values\_from =}\NormalTok{ mean, }\AttributeTok{names\_prefix =} \StringTok{"mean\_"}\NormalTok{)}

\CommentTok{\# Ensure correct column names (assuming \textasciigrave{}+\textasciigrave{} and \textasciigrave{}{-}\textasciigrave{} are original bx values)}
\NormalTok{shaded\_data\_IP\_Hue }\OtherTok{\textless{}{-}}\NormalTok{ shaded\_data\_IP\_Hue }\SpecialCharTok{\%\textgreater{}\%}
  \FunctionTok{rename\_with}\NormalTok{(}\SpecialCharTok{\textasciitilde{}} \FunctionTok{gsub}\NormalTok{(}\StringTok{"mean\_}\SpecialCharTok{\textbackslash{}\textbackslash{}}\StringTok{+"}\NormalTok{, }\StringTok{"mean\_group1"}\NormalTok{, .x)) }\SpecialCharTok{\%\textgreater{}\%}
  \FunctionTok{rename\_with}\NormalTok{(}\SpecialCharTok{\textasciitilde{}} \FunctionTok{gsub}\NormalTok{(}\StringTok{"mean\_}\SpecialCharTok{\textbackslash{}\textbackslash{}}\StringTok{{-}"}\NormalTok{, }\StringTok{"mean\_group2"}\NormalTok{, .x)) }\SpecialCharTok{\%\textgreater{}\%}
  \FunctionTok{mutate}\NormalTok{(}
    \AttributeTok{lower =} \FunctionTok{pmin}\NormalTok{(mean\_group1, mean\_group2, }\AttributeTok{na.rm =} \ConstantTok{TRUE}\NormalTok{),}
    \AttributeTok{upper =} \FunctionTok{pmax}\NormalTok{(mean\_group1, mean\_group2, }\AttributeTok{na.rm =} \ConstantTok{TRUE}\NormalTok{)}
\NormalTok{  )}

\CommentTok{\# Define x{-}axis breaks and labels}
\NormalTok{x\_breaks }\OtherTok{\textless{}{-}} \FunctionTok{seq}\NormalTok{(}\FunctionTok{floor}\NormalTok{(}\FunctionTok{min}\NormalTok{(plant\_region\_areas\_IPHue3}\SpecialCharTok{$}\NormalTok{dag)), }
                \FunctionTok{ceiling}\NormalTok{(}\FunctionTok{max}\NormalTok{(plant\_region\_areas\_IPHue3}\SpecialCharTok{$}\NormalTok{dag)), }\AttributeTok{by =} \DecValTok{5}\NormalTok{)}
\NormalTok{x\_labels }\OtherTok{\textless{}{-}} \FunctionTok{as.character}\NormalTok{(x\_breaks)}

\CommentTok{\# Display final data structure}
\CommentTok{\#head(shaded\_data\_IP\_Hue)}
\end{Highlighting}
\end{Shaded}

Plot the growth for IP-Hue-3 over time

\begin{Shaded}
\begin{Highlighting}[]
\CommentTok{\# Create the plot}
\NormalTok{IP\_Hue\_3 }\OtherTok{\textless{}{-}} \FunctionTok{ggplot}\NormalTok{() }\SpecialCharTok{+}
  \CommentTok{\# Add ribbon without legend}
  \FunctionTok{geom\_ribbon}\NormalTok{(}
    \AttributeTok{data =}\NormalTok{ shaded\_data\_IP\_Hue,}
    \FunctionTok{aes}\NormalTok{(}\AttributeTok{x =}\NormalTok{ dag, }\AttributeTok{ymin =}\NormalTok{ lower, }\AttributeTok{ymax =}\NormalTok{ upper),}
    \AttributeTok{fill =} \StringTok{"\#0072B2"}\NormalTok{, }\AttributeTok{alpha =} \FloatTok{0.2}\NormalTok{,}
    \AttributeTok{show.legend =} \ConstantTok{FALSE}
\NormalTok{  ) }\SpecialCharTok{+}
  \CommentTok{\# Add points with legend}
  \FunctionTok{geom\_point}\NormalTok{(}
    \AttributeTok{data =}\NormalTok{ plant\_region\_areas\_IPHue3,}
    \FunctionTok{aes}\NormalTok{(}\AttributeTok{x =}\NormalTok{ dag, }\AttributeTok{y =}\NormalTok{ mean, }\AttributeTok{colour =}\NormalTok{ bx, }\AttributeTok{fill =}\NormalTok{ bx),}
    \AttributeTok{shape =} \DecValTok{21}\NormalTok{, }\AttributeTok{size =} \FloatTok{2.5}\NormalTok{, }\AttributeTok{alpha =} \FloatTok{0.4}
\NormalTok{  ) }\SpecialCharTok{+}
  \CommentTok{\# Add line without legend}
  \FunctionTok{geom\_line}\NormalTok{(}
    \AttributeTok{data =}\NormalTok{ plant\_region\_areas\_IPHue3,}
    \FunctionTok{aes}\NormalTok{(}\AttributeTok{x =}\NormalTok{ dag, }\AttributeTok{y =}\NormalTok{ mean, }\AttributeTok{colour =}\NormalTok{ bx),}
    \AttributeTok{show.legend =} \ConstantTok{FALSE}
\NormalTok{  ) }\SpecialCharTok{+}
  \CommentTok{\# Modify the y{-}axis label and scale}
  \FunctionTok{ggtitle}\NormalTok{(}\StringTok{"Daily rosette area: IP{-}Hue{-}3"}\NormalTok{) }\SpecialCharTok{+}
  \FunctionTok{labs}\NormalTok{(}
    \AttributeTok{y =} \FunctionTok{expression}\NormalTok{(}\StringTok{"Rosette area (10"}\SpecialCharTok{\^{}}\DecValTok{4}\SpecialCharTok{*}\StringTok{")"}\NormalTok{), }
    \AttributeTok{x =} \StringTok{"Days after germination"}
\NormalTok{  ) }\SpecialCharTok{+}
  \FunctionTok{scale\_y\_continuous}\NormalTok{(}\AttributeTok{labels =} \ControlFlowTok{function}\NormalTok{(x) x }\SpecialCharTok{/} \DecValTok{10000}\NormalTok{) }\SpecialCharTok{+}
  \FunctionTok{scale\_x\_continuous}\NormalTok{(}\AttributeTok{breaks =}\NormalTok{ x\_breaks, }\AttributeTok{labels =}\NormalTok{ x\_labels) }\SpecialCharTok{+}
  \CommentTok{\# Define manual color and fill scales with legend settings}
  \FunctionTok{scale\_colour\_manual}\NormalTok{(}
    \AttributeTok{values =} \FunctionTok{c}\NormalTok{(}\StringTok{"seagreen"}\NormalTok{, }\StringTok{"gold2"}\NormalTok{),}
    \AttributeTok{guide =} \StringTok{"none"} \CommentTok{\# Suppress colour legend}
\NormalTok{  ) }\SpecialCharTok{+}
  \FunctionTok{scale\_fill\_manual}\NormalTok{(}
    \AttributeTok{values =} \FunctionTok{c}\NormalTok{(}\StringTok{"seagreen"}\NormalTok{, }\StringTok{"gold2"}\NormalTok{),}
    \AttributeTok{name =} \StringTok{"Microbiome"}\NormalTok{,  }\CommentTok{\# Update legend title}
    \AttributeTok{labels =} \FunctionTok{c}\NormalTok{(}\StringTok{"BX{-}"}\NormalTok{, }\StringTok{"BX+"}\NormalTok{),  }\CommentTok{\# Update legend labels}
    \AttributeTok{guide =} \FunctionTok{guide\_legend}\NormalTok{(}
      \AttributeTok{override.aes =} \FunctionTok{list}\NormalTok{(}\AttributeTok{shape =} \DecValTok{21}\NormalTok{, }\AttributeTok{size =} \FloatTok{2.5}\NormalTok{, }\AttributeTok{alpha =} \FloatTok{0.4}\NormalTok{) }\CommentTok{\# Match point aesthetics}
\NormalTok{    )}
\NormalTok{  ) }\SpecialCharTok{+}
  \FunctionTok{theme\_bw}\NormalTok{() }\SpecialCharTok{+}
  \FunctionTok{theme}\NormalTok{(}
    \AttributeTok{axis.text =} \FunctionTok{element\_text}\NormalTok{(}\AttributeTok{size =} \DecValTok{7}\NormalTok{),}
    \AttributeTok{legend.title =} \FunctionTok{element\_text}\NormalTok{(}\AttributeTok{size =} \DecValTok{7}\NormalTok{),}
    \AttributeTok{legend.text =} \FunctionTok{element\_text}\NormalTok{(}\AttributeTok{size =} \DecValTok{7}\NormalTok{),}
    \AttributeTok{plot.title =} \FunctionTok{element\_text}\NormalTok{(}\AttributeTok{size =} \DecValTok{7}\NormalTok{),}
    \AttributeTok{axis.title.y =} \FunctionTok{element\_text}\NormalTok{(}\AttributeTok{size =} \DecValTok{7}\NormalTok{),}
    \AttributeTok{axis.title.x =} \FunctionTok{element\_text}\NormalTok{(}\AttributeTok{size =} \DecValTok{7}\NormalTok{),}
    \AttributeTok{legend.position =} \FunctionTok{c}\NormalTok{(}\FloatTok{0.2}\NormalTok{, }\FloatTok{0.95}\NormalTok{),  }\CommentTok{\# Position legend in top{-}right}
    \AttributeTok{legend.justification =} \FunctionTok{c}\NormalTok{(}\DecValTok{1}\NormalTok{, }\DecValTok{1}\NormalTok{),   }\CommentTok{\# Align legend box to top{-}right}
    \AttributeTok{legend.background =} \FunctionTok{element\_rect}\NormalTok{(}\AttributeTok{fill =} \StringTok{"white"}\NormalTok{, }\AttributeTok{color =} \StringTok{"black"}\NormalTok{)}
\NormalTok{  )}

\CommentTok{\# Display the plot}
\FunctionTok{print}\NormalTok{(IP\_Hue\_3)}
\end{Highlighting}
\end{Shaded}

\begin{Shaded}
\begin{Highlighting}[]
\CommentTok{\# Purpose: Displays the plot.}

\CommentTok{\# Save the plot}
\FunctionTok{ggsave}\NormalTok{(}\StringTok{"IP\_hue\_3\_growth.pdf"}\NormalTok{, }\AttributeTok{plot =}\NormalTok{ IP\_Hue\_3, }\AttributeTok{width =} \DecValTok{10}\NormalTok{, }\AttributeTok{height =} \DecValTok{10}\NormalTok{, }\AttributeTok{units =} \StringTok{"cm"}\NormalTok{, }\AttributeTok{path =} \StringTok{"Output/Figures/"}\NormalTok{)}
\CommentTok{\# Purpose: Saves the plot as a pdf file.}
\end{Highlighting}
\end{Shaded}

Plot the growth feedback for IP-Hue-3 over time

\textbf{Purpose}: Calculates the difference in mean plant region area
between the two treatments for each day and creates a bar plot showing
the \emph{daily growth feedback} for IP-Hue-3.

\begin{Shaded}
\begin{Highlighting}[]
\CommentTok{\# Compute the difference between the two groups}
\NormalTok{shaded\_data\_IP\_Hue }\OtherTok{\textless{}{-}}\NormalTok{ shaded\_data\_IP\_Hue }\SpecialCharTok{\%\textgreater{}\%}
  \FunctionTok{mutate}\NormalTok{(}\AttributeTok{difference =}\NormalTok{ mean\_group1 }\SpecialCharTok{{-}}\NormalTok{ mean\_group2)}
\CommentTok{\# Purpose: Calculates the difference in mean plant region area between the two treatments for each day.}

\CommentTok{\# Plot with shaded area and differences as bars}
\NormalTok{IP\_Hue\_3\_response }\OtherTok{\textless{}{-}} \FunctionTok{ggplot}\NormalTok{() }\SpecialCharTok{+}
  \FunctionTok{geom\_bar}\NormalTok{(}\AttributeTok{data =}\NormalTok{ shaded\_data\_IP\_Hue, }\FunctionTok{aes}\NormalTok{(}\AttributeTok{x =}\NormalTok{ dag, }\AttributeTok{y =}\NormalTok{ difference), }\AttributeTok{stat =} \StringTok{"identity"}\NormalTok{, }\AttributeTok{fill =} \StringTok{"\#0072B2"}\NormalTok{, }\AttributeTok{alpha =} \FloatTok{0.8}\NormalTok{) }\SpecialCharTok{+}
  \FunctionTok{ggtitle}\NormalTok{(}\StringTok{"Daily growth feeedbacks: IP{-}Hue{-}3"}\NormalTok{) }\SpecialCharTok{+}
  \FunctionTok{labs}\NormalTok{(}
    \AttributeTok{y =} \FunctionTok{expression}\NormalTok{(}\StringTok{"Growth feedback (10"}\SpecialCharTok{\^{}}\DecValTok{4}\SpecialCharTok{*}\StringTok{")"}\NormalTok{), }
    \AttributeTok{x =} \StringTok{"Days after germination"}
\NormalTok{  ) }\SpecialCharTok{+}
  \FunctionTok{scale\_y\_continuous}\NormalTok{(}\AttributeTok{labels =} \ControlFlowTok{function}\NormalTok{(x) x }\SpecialCharTok{/} \DecValTok{10000}\NormalTok{) }\SpecialCharTok{+}
  \FunctionTok{scale\_x\_continuous}\NormalTok{(}\AttributeTok{breaks =}\NormalTok{ x\_breaks, }\AttributeTok{labels =}\NormalTok{ x\_labels) }\SpecialCharTok{+}
  \FunctionTok{theme\_bw}\NormalTok{() }\SpecialCharTok{+}
  \FunctionTok{theme}\NormalTok{(}\AttributeTok{axis.text =} \FunctionTok{element\_text}\NormalTok{(}\AttributeTok{size =} \DecValTok{7}\NormalTok{),}
        \AttributeTok{legend.title =} \FunctionTok{element\_text}\NormalTok{(}\AttributeTok{size =} \DecValTok{7}\NormalTok{),}
        \AttributeTok{legend.text =} \FunctionTok{element\_text}\NormalTok{(}\AttributeTok{size =} \DecValTok{7}\NormalTok{),}
        \AttributeTok{plot.title =} \FunctionTok{element\_text}\NormalTok{(}\AttributeTok{size =} \DecValTok{7}\NormalTok{),}
        \AttributeTok{axis.title.y =} \FunctionTok{element\_text}\NormalTok{(}\AttributeTok{size =} \DecValTok{7}\NormalTok{),}
        \AttributeTok{axis.title.x =} \FunctionTok{element\_text}\NormalTok{(}\AttributeTok{size =} \DecValTok{7}\NormalTok{),}
        \AttributeTok{legend.position =} \StringTok{"none"}\NormalTok{)       }\CommentTok{\# Remove legend}
\CommentTok{\# Purpose: Creates a bar plot showing the daily growth feedback for IP{-}Hue{-}3, calculated as the difference between the two treatments.}

\CommentTok{\# Display the plot}
\FunctionTok{print}\NormalTok{(IP\_Hue\_3\_response)}
\end{Highlighting}
\end{Shaded}

\begin{Shaded}
\begin{Highlighting}[]
\CommentTok{\# Purpose: Displays the plot.}

\CommentTok{\# Save the plot}
\FunctionTok{ggsave}\NormalTok{(}\StringTok{"IP\_Hue\_3\_response.pdf"}\NormalTok{, }\AttributeTok{plot =}\NormalTok{ IP\_Hue\_3\_response, }\AttributeTok{width =} \DecValTok{10}\NormalTok{, }\AttributeTok{height =} \DecValTok{10}\NormalTok{, }\AttributeTok{units =} \StringTok{"cm"}\NormalTok{, }\AttributeTok{path =} \StringTok{"Output/Figures/"}\NormalTok{)}
\CommentTok{\# Purpose: Saves the plot as a PNG file.}
\end{Highlighting}
\end{Shaded}

\subsubsection{Hadd-3}\label{hadd-3}

Filter for Hadd-3 and Plot the Growth Over Time

\textbf{Purpose}: Filters the data to include only the genotype with
genotype\_id 9392 (Hadd-3), calculates the mean plant region area for
each combination of bx and dag, removes duplicates and irrelevant
columns, pivots the data to spread `bx' treatment values into separate
columns, calculates the area difference, and ranks the genotypes by the
area difference.

\begin{Shaded}
\begin{Highlighting}[]
\CommentTok{\# Filter data for genotype\_id 9392}
\NormalTok{plant\_region\_areas\_Hadd3 }\OtherTok{\textless{}{-}}\NormalTok{ plant\_region\_areas }\SpecialCharTok{\%\textgreater{}\%} 
  \FunctionTok{filter}\NormalTok{(dag }\SpecialCharTok{\textless{}} \DecValTok{21}\NormalTok{, genotype\_id }\SpecialCharTok{==} \DecValTok{9392}\NormalTok{)}
\CommentTok{\# Purpose: Filters the dataset to include only data for the accession with genotype\_id 9392 (Hadd{-}3) and days less than 20.}

\CommentTok{\# Calculate mean plant region area for each combination of bx and dag}
\NormalTok{plant\_region\_areas\_Hadd3 }\OtherTok{\textless{}{-}}\NormalTok{ plant\_region\_areas\_Hadd3 }\SpecialCharTok{\%\textgreater{}\%}
  \FunctionTok{group\_by}\NormalTok{(bx, dag) }\SpecialCharTok{\%\textgreater{}\%}
  \FunctionTok{mutate}\NormalTok{(}\AttributeTok{mean =} \FunctionTok{mean}\NormalTok{(plant\_region\_area, }\AttributeTok{na.rm =} \ConstantTok{TRUE}\NormalTok{)) }\SpecialCharTok{\%\textgreater{}\%}
  \FunctionTok{ungroup}\NormalTok{()}
\CommentTok{\# Purpose: Calculates the mean plant region area for each treatment (bx) and day (dag) combination.}

\CommentTok{\# Remove unnecessary columns and keep only unique entries}
\NormalTok{plant\_region\_areas\_Hadd3 }\OtherTok{\textless{}{-}}\NormalTok{ plant\_region\_areas\_Hadd3 }\SpecialCharTok{\%\textgreater{}\%}
  \FunctionTok{select}\NormalTok{(}\SpecialCharTok{{-}}\FunctionTok{c}\NormalTok{(}\DecValTok{1}\NormalTok{, }\DecValTok{6}\NormalTok{)) }\SpecialCharTok{\%\textgreater{}\%}
  \FunctionTok{distinct}\NormalTok{()}
\CommentTok{\# Purpose: Removes unnecessary columns and ensures that only unique rows remain in the dataset.}

\CommentTok{\# Create a new data frame for the shaded area between curves}
\CommentTok{\# Spread the data to have separate columns for each group}
\NormalTok{shaded\_data\_Hadd3 }\OtherTok{\textless{}{-}}\NormalTok{ plant\_region\_areas\_Hadd3 }\SpecialCharTok{\%\textgreater{}\%}
  \FunctionTok{pivot\_wider}\NormalTok{(}\AttributeTok{names\_from =}\NormalTok{ bx, }\AttributeTok{values\_from =}\NormalTok{ mean)}
\CommentTok{\# Purpose: Transforms the data to have separate columns for each treatment, preparing it for shaded area plotting.}

\CommentTok{\# Ensure the correct names for the new columns}
\CommentTok{\# Assuming the groups are "group1" and "group2"}
\NormalTok{shaded\_data\_Hadd3 }\OtherTok{\textless{}{-}}\NormalTok{ shaded\_data\_Hadd3 }\SpecialCharTok{\%\textgreater{}\%}
  \FunctionTok{rename}\NormalTok{(}\AttributeTok{mean\_group1 =} \StringTok{\textasciigrave{}}\AttributeTok{+}\StringTok{\textasciigrave{}}\NormalTok{, }\AttributeTok{mean\_group2 =} \StringTok{\textasciigrave{}}\AttributeTok{{-}}\StringTok{\textasciigrave{}}\NormalTok{) }\SpecialCharTok{\%\textgreater{}\%}
  \FunctionTok{mutate}\NormalTok{(}\AttributeTok{lower =} \FunctionTok{pmin}\NormalTok{(mean\_group1, mean\_group2, }\AttributeTok{na.rm =} \ConstantTok{TRUE}\NormalTok{),}
         \AttributeTok{upper =} \FunctionTok{pmax}\NormalTok{(mean\_group1, mean\_group2, }\AttributeTok{na.rm =} \ConstantTok{TRUE}\NormalTok{))}
\CommentTok{\# Purpose: Renames the columns for clarity and calculates the lower and upper bounds for the shaded area.}
\end{Highlighting}
\end{Shaded}

Plot the growth for Hadd-3 over time

\begin{Shaded}
\begin{Highlighting}[]
\CommentTok{\# Create the plot}
\NormalTok{Hadd\_3 }\OtherTok{\textless{}{-}} \FunctionTok{ggplot}\NormalTok{() }\SpecialCharTok{+}
  \CommentTok{\# Add ribbon without legend}
  \FunctionTok{geom\_ribbon}\NormalTok{(}
    \AttributeTok{data =}\NormalTok{ shaded\_data\_Hadd3,}
    \FunctionTok{aes}\NormalTok{(}\AttributeTok{x =}\NormalTok{ dag, }\AttributeTok{ymin =}\NormalTok{ lower, }\AttributeTok{ymax =}\NormalTok{ upper),}
    \AttributeTok{fill =} \StringTok{"orange2"}\NormalTok{, }\AttributeTok{alpha =} \FloatTok{0.2}\NormalTok{,}
    \AttributeTok{show.legend =} \ConstantTok{FALSE}
\NormalTok{  ) }\SpecialCharTok{+}
  \CommentTok{\# Add points with legend}
  \FunctionTok{geom\_point}\NormalTok{(}
    \AttributeTok{data =}\NormalTok{ plant\_region\_areas\_Hadd3,}
    \FunctionTok{aes}\NormalTok{(}\AttributeTok{x =}\NormalTok{ dag, }\AttributeTok{y =}\NormalTok{ mean, }\AttributeTok{colour =}\NormalTok{ bx, }\AttributeTok{fill =}\NormalTok{ bx),}
    \AttributeTok{shape =} \DecValTok{21}\NormalTok{, }\AttributeTok{size =} \FloatTok{2.5}\NormalTok{, }\AttributeTok{alpha =} \FloatTok{0.4}
\NormalTok{  ) }\SpecialCharTok{+}
  \CommentTok{\# Add line without legend}
  \FunctionTok{geom\_line}\NormalTok{(}
    \AttributeTok{data =}\NormalTok{ plant\_region\_areas\_Hadd3,}
    \FunctionTok{aes}\NormalTok{(}\AttributeTok{x =}\NormalTok{ dag, }\AttributeTok{y =}\NormalTok{ mean, }\AttributeTok{colour =}\NormalTok{ bx),}
    \AttributeTok{show.legend =} \ConstantTok{FALSE}
\NormalTok{  ) }\SpecialCharTok{+}
  \CommentTok{\# Title and axis labels}
  \FunctionTok{ggtitle}\NormalTok{(}\StringTok{"Daily rosette area: Hadd{-}3"}\NormalTok{) }\SpecialCharTok{+}
  \FunctionTok{labs}\NormalTok{(}
    \AttributeTok{y =} \FunctionTok{expression}\NormalTok{(}\StringTok{"Rosette area (10"}\SpecialCharTok{\^{}}\DecValTok{4}\SpecialCharTok{*}\StringTok{")"}\NormalTok{), }
    \AttributeTok{x =} \StringTok{"Days after germination"}
\NormalTok{  ) }\SpecialCharTok{+}
  \CommentTok{\# Scales for y{-}axis and x{-}axis}
  \FunctionTok{scale\_y\_continuous}\NormalTok{(}\AttributeTok{labels =} \ControlFlowTok{function}\NormalTok{(x) x }\SpecialCharTok{/} \DecValTok{10000}\NormalTok{) }\SpecialCharTok{+}
  \FunctionTok{scale\_x\_continuous}\NormalTok{(}\AttributeTok{breaks =}\NormalTok{ x\_breaks, }\AttributeTok{labels =}\NormalTok{ x\_labels) }\SpecialCharTok{+}
  \CommentTok{\# Define manual color and fill scales with legend settings}
  \FunctionTok{scale\_colour\_manual}\NormalTok{(}
    \AttributeTok{values =} \FunctionTok{c}\NormalTok{(}\StringTok{"seagreen"}\NormalTok{, }\StringTok{"gold2"}\NormalTok{),}
    \AttributeTok{guide =} \StringTok{"none"} \CommentTok{\# Suppress colour legend}
\NormalTok{  ) }\SpecialCharTok{+}
  \FunctionTok{scale\_fill\_manual}\NormalTok{(}
    \AttributeTok{values =} \FunctionTok{c}\NormalTok{(}\StringTok{"seagreen"}\NormalTok{, }\StringTok{"gold2"}\NormalTok{),}
    \AttributeTok{name =} \StringTok{"Microbiome"}\NormalTok{,  }\CommentTok{\# Update legend title}
    \AttributeTok{labels =} \FunctionTok{c}\NormalTok{(}\StringTok{"BX{-}"}\NormalTok{, }\StringTok{"BX+"}\NormalTok{),  }\CommentTok{\# Update legend labels}
    \AttributeTok{guide =} \FunctionTok{guide\_legend}\NormalTok{(}
      \AttributeTok{override.aes =} \FunctionTok{list}\NormalTok{(}\AttributeTok{shape =} \DecValTok{21}\NormalTok{, }\AttributeTok{size =} \FloatTok{2.5}\NormalTok{, }\AttributeTok{alpha =} \FloatTok{0.4}\NormalTok{) }\CommentTok{\# Match point aesthetics}
\NormalTok{    )}
\NormalTok{  ) }\SpecialCharTok{+}
  \CommentTok{\# Theme adjustments}
  \FunctionTok{theme\_bw}\NormalTok{() }\SpecialCharTok{+}
  \FunctionTok{theme}\NormalTok{(}
    \AttributeTok{axis.text =} \FunctionTok{element\_text}\NormalTok{(}\AttributeTok{size =} \DecValTok{7}\NormalTok{),}
    \AttributeTok{legend.title =} \FunctionTok{element\_text}\NormalTok{(}\AttributeTok{size =} \DecValTok{7}\NormalTok{),}
    \AttributeTok{legend.text =} \FunctionTok{element\_text}\NormalTok{(}\AttributeTok{size =} \DecValTok{7}\NormalTok{),}
    \AttributeTok{plot.title =} \FunctionTok{element\_text}\NormalTok{(}\AttributeTok{size =} \DecValTok{7}\NormalTok{),}
    \AttributeTok{axis.title.y =} \FunctionTok{element\_text}\NormalTok{(}\AttributeTok{size =} \DecValTok{7}\NormalTok{),}
    \AttributeTok{axis.title.x =} \FunctionTok{element\_text}\NormalTok{(}\AttributeTok{size =} \DecValTok{7}\NormalTok{),}
    \AttributeTok{legend.position =} \FunctionTok{c}\NormalTok{(}\FloatTok{0.2}\NormalTok{, }\FloatTok{0.95}\NormalTok{),  }\CommentTok{\# Position legend in top{-}right}
    \AttributeTok{legend.justification =} \FunctionTok{c}\NormalTok{(}\DecValTok{1}\NormalTok{, }\DecValTok{1}\NormalTok{),   }\CommentTok{\# Align legend box to top{-}right}
    \AttributeTok{legend.background =} \FunctionTok{element\_rect}\NormalTok{(}\AttributeTok{fill =} \StringTok{"white"}\NormalTok{, }\AttributeTok{color =} \StringTok{"black"}\NormalTok{)}
\NormalTok{  )}

\CommentTok{\# Display the plot}
\NormalTok{Hadd\_3}
\end{Highlighting}
\end{Shaded}

\begin{Shaded}
\begin{Highlighting}[]
\FunctionTok{ggsave}\NormalTok{(}\StringTok{"Hadd\_3\_growth.pdf"}\NormalTok{, }\AttributeTok{plot =}\NormalTok{ Hadd\_3, }\AttributeTok{width =} \DecValTok{10}\NormalTok{, }\AttributeTok{height =} \DecValTok{10}\NormalTok{, }\AttributeTok{units =} \StringTok{"cm"}\NormalTok{, }\AttributeTok{path =} \StringTok{"Output/Figures/"}\NormalTok{)}
\end{Highlighting}
\end{Shaded}

Plot the growth feedback for Hadd-3 over time

\textbf{Purpose}: Calculates the difference in mean plant region area
between the two treatments for each day and creates a bar plot showing
the \emph{daily growth feedback} for Hadd-3.

\begin{Shaded}
\begin{Highlighting}[]
\CommentTok{\# Compute the difference between the two groups}
\NormalTok{shaded\_data\_Hadd3 }\OtherTok{\textless{}{-}}\NormalTok{ shaded\_data\_Hadd3 }\SpecialCharTok{\%\textgreater{}\%}
  \FunctionTok{mutate}\NormalTok{(}\AttributeTok{difference =}\NormalTok{ mean\_group1 }\SpecialCharTok{{-}}\NormalTok{ mean\_group2)}
\CommentTok{\# Purpose: Calculates the difference in mean plant region area between the two treatments for each day.}

\CommentTok{\# Plot with shaded area and differences as bars}
\NormalTok{Hadd3\_response }\OtherTok{\textless{}{-}} \FunctionTok{ggplot}\NormalTok{() }\SpecialCharTok{+}
  \FunctionTok{geom\_bar}\NormalTok{(}\AttributeTok{data =}\NormalTok{ shaded\_data\_Hadd3, }\FunctionTok{aes}\NormalTok{(}\AttributeTok{x =}\NormalTok{ dag, }\AttributeTok{y =}\NormalTok{ difference), }\AttributeTok{stat =} \StringTok{"identity"}\NormalTok{, }\AttributeTok{fill =} \StringTok{"orange2"}\NormalTok{, }\AttributeTok{alpha =} \FloatTok{0.8}\NormalTok{) }\SpecialCharTok{+}
  \FunctionTok{ggtitle}\NormalTok{(}\StringTok{"Daily growth feeedbacks: Hadd{-}3"}\NormalTok{) }\SpecialCharTok{+}
  \FunctionTok{labs}\NormalTok{(}
    \AttributeTok{y =} \FunctionTok{expression}\NormalTok{(}\StringTok{"Growth feedback (10"}\SpecialCharTok{\^{}}\DecValTok{4}\SpecialCharTok{*}\StringTok{")"}\NormalTok{), }
    \AttributeTok{x =} \StringTok{"Days after germination"}
\NormalTok{  ) }\SpecialCharTok{+}
  \FunctionTok{scale\_y\_continuous}\NormalTok{(}\AttributeTok{labels =} \ControlFlowTok{function}\NormalTok{(x) x }\SpecialCharTok{/} \DecValTok{10000}\NormalTok{) }\SpecialCharTok{+}
  \FunctionTok{scale\_x\_continuous}\NormalTok{(}\AttributeTok{breaks =}\NormalTok{ x\_breaks, }\AttributeTok{labels =}\NormalTok{ x\_labels) }\SpecialCharTok{+}
  \FunctionTok{theme\_bw}\NormalTok{() }\SpecialCharTok{+}
  \FunctionTok{theme}\NormalTok{(}\AttributeTok{axis.text =} \FunctionTok{element\_text}\NormalTok{(}\AttributeTok{size =} \DecValTok{7}\NormalTok{),}
        \AttributeTok{legend.title =} \FunctionTok{element\_text}\NormalTok{(}\AttributeTok{size =} \DecValTok{7}\NormalTok{),}
        \AttributeTok{legend.text =} \FunctionTok{element\_text}\NormalTok{(}\AttributeTok{size =} \DecValTok{7}\NormalTok{),}
        \AttributeTok{plot.title =} \FunctionTok{element\_text}\NormalTok{(}\AttributeTok{size =} \DecValTok{7}\NormalTok{),}
        \AttributeTok{axis.title.y =} \FunctionTok{element\_text}\NormalTok{(}\AttributeTok{size =} \DecValTok{7}\NormalTok{),}
        \AttributeTok{axis.title.x =} \FunctionTok{element\_text}\NormalTok{(}\AttributeTok{size =} \DecValTok{7}\NormalTok{),}
        \CommentTok{\#axis.text.y = element\_blank(),  \# Remove y{-}axis text}
        \AttributeTok{legend.position =} \StringTok{"none"}\NormalTok{)       }\CommentTok{\# Remove legend}
\CommentTok{\# Purpose: Creates a bar plot showing the daily growth feedback for Hadd{-}3, calculated as the difference between the two treatments.}

\CommentTok{\# Display the plot}
\FunctionTok{print}\NormalTok{(Hadd3\_response)}
\end{Highlighting}
\end{Shaded}

\begin{Shaded}
\begin{Highlighting}[]
\CommentTok{\# Purpose: Displays the plot.}

\CommentTok{\# Save the plot}
\FunctionTok{ggsave}\NormalTok{(}\StringTok{"Hadd\_3\_response.pdf"}\NormalTok{, }\AttributeTok{plot =}\NormalTok{ Hadd3\_response, }\AttributeTok{width =} \DecValTok{10}\NormalTok{, }\AttributeTok{height =} \DecValTok{10}\NormalTok{, }\AttributeTok{units =} \StringTok{"cm"}\NormalTok{, }\AttributeTok{path =} \StringTok{"Output/Figures/"}\NormalTok{)}
\end{Highlighting}
\end{Shaded}

\begin{figure}[H]

{\centering \pandocbounded{\includegraphics[keepaspectratio]{01_Accession_phenotypes_files/figure-pdf/Fig1b-Growth-feedback-extremeACC-1.pdf}}

}

\caption{Combined plots showing the growth and growth feedbacks of the
two most exctreme accessions.}

\end{figure}%

\section{Geographic distribution}\label{geographic-distribution}

\subsection{Import data}\label{import-data}

\begin{Shaded}
\begin{Highlighting}[]
\NormalTok{ACC\_coordinates }\OtherTok{\textless{}{-}} \FunctionTok{read\_excel}\NormalTok{(}\StringTok{"Input/Dataset1.xlsx"}\NormalTok{, }\AttributeTok{sheet =} \StringTok{"accession\_coordinates"}\NormalTok{)}
\FunctionTok{kable}\NormalTok{(ACC\_coordinates[}\DecValTok{1}\SpecialCharTok{:}\DecValTok{3}\NormalTok{, ], }\AttributeTok{caption =} \StringTok{"Coordinate table."}\NormalTok{)}
\end{Highlighting}
\end{Shaded}

\begin{longtable}[]{@{}rllrr@{}}
\caption{Coordinate table.}\tabularnewline
\toprule\noalign{}
ACC\_ID & Country & Region & Lat & Long \\
\midrule\noalign{}
\endfirsthead
\toprule\noalign{}
ACC\_ID & Country & Region & Lat & Long \\
\midrule\noalign{}
\endhead
\bottomrule\noalign{}
\endlastfoot
5921 & CZE & DraIV 3 & 49.4112 & 16.2815 \\
5867 & SWE & Dra2 & 55.7600 & 14.1200 \\
10006 & UKR & Kastel Mountain & 44.6419 & 34.3814 \\
\end{longtable}

\subsection{Correlation of growth feedback and geographical
location}\label{correlation-of-growth-feedback-and-geographical-location}

\subsubsection{Prepare data for correlation and regression
analysis}\label{prepare-data-for-correlation-and-regression-analysis}

\textbf{Purpose:} Filters the data to include only observations from day
10, assigns a preference based on the area difference, sorts by area
difference, and ranks the accessions.

\begin{Shaded}
\begin{Highlighting}[]
\CommentTok{\# Filter for observations on day 10, assign preference, and sort}
\NormalTok{plant\_region\_areas\_DAG10 }\OtherTok{\textless{}{-}}\NormalTok{ plant\_region\_areas\_DAG20 }\SpecialCharTok{\%\textgreater{}\%}
  \FunctionTok{filter}\NormalTok{(dag }\SpecialCharTok{==} \DecValTok{15}\NormalTok{) }\SpecialCharTok{\%\textgreater{}\%}
  \FunctionTok{mutate}\NormalTok{(}\AttributeTok{preference =} \FunctionTok{case\_when}\NormalTok{(}
\NormalTok{    area\_diff }\SpecialCharTok{\textgreater{}=} \DecValTok{0} \SpecialCharTok{\textasciitilde{}} \StringTok{"plus"}\NormalTok{,}
    \ConstantTok{TRUE} \SpecialCharTok{\textasciitilde{}} \StringTok{"minus"}
\NormalTok{  )) }\SpecialCharTok{\%\textgreater{}\%}
  \FunctionTok{arrange}\NormalTok{(}\FunctionTok{desc}\NormalTok{(area\_diff)) }\SpecialCharTok{\%\textgreater{}\%}
  \FunctionTok{mutate}\NormalTok{(}\AttributeTok{rank =} \FunctionTok{row\_number}\NormalTok{())}
\end{Highlighting}
\end{Shaded}

Merge the data tables

\textbf{Purpose:} Merges the phenotype data
(plant\_region\_areas\_DAG10) with the coordinates data
(ACC\_coordinates) based on genotype IDs.

\begin{Shaded}
\begin{Highlighting}[]
\CommentTok{\# Merge the data tables and include only Lat and Long columns from ACC\_coordinates}
\NormalTok{coordinates\_merged\_data }\OtherTok{\textless{}{-}}\NormalTok{ plant\_region\_areas\_DAG10 }\SpecialCharTok{\%\textgreater{}\%}
  \FunctionTok{left\_join}\NormalTok{(}
\NormalTok{    ACC\_coordinates }\SpecialCharTok{\%\textgreater{}\%} \FunctionTok{select}\NormalTok{(ACC\_ID, Lat, Long),  }\CommentTok{\# Select only relevant columns}
    \AttributeTok{by =} \FunctionTok{c}\NormalTok{(}\StringTok{"genotype\_id"} \OtherTok{=} \StringTok{"ACC\_ID"}\NormalTok{)               }\CommentTok{\# Specify the join key}
\NormalTok{  )}
\end{Highlighting}
\end{Shaded}

\subsubsection{\texorpdfstring{\textbf{Pearson correlation} and
Regression
Analysis}{Pearson correlation and Regression Analysis}}\label{pearson-correlation-and-regression-analysis}

\begin{Shaded}
\begin{Highlighting}[]
\CommentTok{\# Perform Pearson correlation test for Latitude}
\NormalTok{cor\_test\_lat }\OtherTok{\textless{}{-}} \FunctionTok{cor.test}\NormalTok{(coordinates\_merged\_data}\SpecialCharTok{$}\NormalTok{Lat, coordinates\_merged\_data}\SpecialCharTok{$}\NormalTok{area\_diff, }
                         \AttributeTok{method =} \StringTok{"pearson"}\NormalTok{, }\AttributeTok{use =} \StringTok{"complete.obs"}\NormalTok{)}

\CommentTok{\# Perform Pearson correlation test for Longitude}
\NormalTok{cor\_test\_long }\OtherTok{\textless{}{-}} \FunctionTok{cor.test}\NormalTok{(coordinates\_merged\_data}\SpecialCharTok{$}\NormalTok{Long, coordinates\_merged\_data}\SpecialCharTok{$}\NormalTok{area\_diff, }
                          \AttributeTok{method =} \StringTok{"pearson"}\NormalTok{, }\AttributeTok{use =} \StringTok{"complete.obs"}\NormalTok{)}

\CommentTok{\# Function to extract correlation test results safely}
\NormalTok{extract\_cor\_results }\OtherTok{\textless{}{-}} \ControlFlowTok{function}\NormalTok{(test\_result, variable\_name) \{}
  \FunctionTok{data.frame}\NormalTok{(}
    \AttributeTok{Variable =}\NormalTok{ variable\_name,}
    \AttributeTok{Correlation =} \FunctionTok{round}\NormalTok{(test\_result}\SpecialCharTok{$}\NormalTok{estimate, }\DecValTok{3}\NormalTok{),  }\CommentTok{\# Rounded for readability}
    \AttributeTok{p\_value =} \FunctionTok{round}\NormalTok{(test\_result}\SpecialCharTok{$}\NormalTok{p.value, }\DecValTok{4}\NormalTok{),}
    \AttributeTok{CI\_lower =} \ControlFlowTok{if}\NormalTok{ (}\SpecialCharTok{!}\FunctionTok{is.null}\NormalTok{(test\_result}\SpecialCharTok{$}\NormalTok{conf.int)) }\FunctionTok{round}\NormalTok{(test\_result}\SpecialCharTok{$}\NormalTok{conf.int[}\DecValTok{1}\NormalTok{], }\DecValTok{3}\NormalTok{) }\ControlFlowTok{else} \ConstantTok{NA}\NormalTok{,}
    \AttributeTok{CI\_upper =} \ControlFlowTok{if}\NormalTok{ (}\SpecialCharTok{!}\FunctionTok{is.null}\NormalTok{(test\_result}\SpecialCharTok{$}\NormalTok{conf.int)) }\FunctionTok{round}\NormalTok{(test\_result}\SpecialCharTok{$}\NormalTok{conf.int[}\DecValTok{2}\NormalTok{], }\DecValTok{3}\NormalTok{) }\ControlFlowTok{else} \ConstantTok{NA}
\NormalTok{  )}
\NormalTok{\}}

\CommentTok{\# Extract results}
\NormalTok{lat\_results }\OtherTok{\textless{}{-}} \FunctionTok{extract\_cor\_results}\NormalTok{(cor\_test\_lat, }\StringTok{"Latitude"}\NormalTok{)}
\NormalTok{long\_results }\OtherTok{\textless{}{-}} \FunctionTok{extract\_cor\_results}\NormalTok{(cor\_test\_long, }\StringTok{"Longitude"}\NormalTok{)}

\CommentTok{\# Combine both results into a single table}
\NormalTok{cor\_results }\OtherTok{\textless{}{-}} \FunctionTok{bind\_rows}\NormalTok{(lat\_results, long\_results)}

\CommentTok{\# Display the results as a neat table}
\FunctionTok{kable}\NormalTok{(cor\_results, }\AttributeTok{caption =} \StringTok{"Correlation Test Results for Latitude and Longitude"}\NormalTok{, }\AttributeTok{digits =} \DecValTok{4}\NormalTok{)}
\end{Highlighting}
\end{Shaded}

\begin{longtable}[]{@{}llrrrr@{}}
\caption{Correlation Test Results for Latitude and
Longitude}\tabularnewline
\toprule\noalign{}
& Variable & Correlation & p\_value & CI\_lower & CI\_upper \\
\midrule\noalign{}
\endfirsthead
\toprule\noalign{}
& Variable & Correlation & p\_value & CI\_lower & CI\_upper \\
\midrule\noalign{}
\endhead
\bottomrule\noalign{}
\endlastfoot
cor\ldots1 & Latitude & 0.002 & 0.9751 & -0.096 & 0.099 \\
cor\ldots2 & Longitude & 0.058 & 0.2398 & -0.039 & 0.155 \\
\end{longtable}

\subsubsection{\texorpdfstring{\textbf{Linear regression
model}}{Linear regression model}}\label{linear-regression-model}

\begin{Shaded}
\begin{Highlighting}[]
\CommentTok{\# Linear regression with Latitude}
\NormalTok{model\_lat }\OtherTok{\textless{}{-}} \FunctionTok{lm}\NormalTok{(area\_diff }\SpecialCharTok{\textasciitilde{}}\NormalTok{ Lat }\SpecialCharTok{+}\NormalTok{ genotype\_id, }\AttributeTok{data =}\NormalTok{ coordinates\_merged\_data)}
\NormalTok{summary\_lat }\OtherTok{\textless{}{-}} \FunctionTok{summary}\NormalTok{(model\_lat)}

\CommentTok{\# Extract coefficients and R{-}squared}
\NormalTok{intercept\_lat }\OtherTok{\textless{}{-}} \FunctionTok{coef}\NormalTok{(summary\_lat)[}\DecValTok{1}\NormalTok{, }\DecValTok{1}\NormalTok{]}
\NormalTok{slope\_lat }\OtherTok{\textless{}{-}} \FunctionTok{coef}\NormalTok{(summary\_lat)[}\DecValTok{2}\NormalTok{, }\DecValTok{1}\NormalTok{]}
\NormalTok{r\_squared\_lat }\OtherTok{\textless{}{-}}\NormalTok{ summary\_lat}\SpecialCharTok{$}\NormalTok{r.squared}
\NormalTok{p\_value\_lat }\OtherTok{\textless{}{-}} \FunctionTok{coef}\NormalTok{(summary\_lat)[}\DecValTok{2}\NormalTok{, }\DecValTok{4}\NormalTok{]}
\NormalTok{p\_value\_f\_lat }\OtherTok{\textless{}{-}}\NormalTok{ summary\_lat}\SpecialCharTok{$}\NormalTok{fstatistic[}\DecValTok{1}\NormalTok{] }\CommentTok{\# F{-}statistic value}
\FunctionTok{print}\NormalTok{(summary\_lat)}
\end{Highlighting}
\end{Shaded}

\begin{verbatim}

Call:
lm(formula = area_diff ~ Lat + genotype_id, data = coordinates_merged_data)

Residuals:
   Min     1Q Median     3Q    Max 
-29305  -4220    489   4132  47752 

Coefficients:
              Estimate Std. Error t value Pr(>|t|)
(Intercept) -948.88265 3679.18872  -0.258    0.797
Lat           -0.86992   62.11497  -0.014    0.989
genotype_id   -0.03179    0.18395  -0.173    0.863

Residual standard error: 8755 on 405 degrees of freedom
Multiple R-squared:  7.612e-05, Adjusted R-squared:  -0.004862 
F-statistic: 0.01542 on 2 and 405 DF,  p-value: 0.9847
\end{verbatim}

\begin{Shaded}
\begin{Highlighting}[]
\ControlFlowTok{if}\NormalTok{ (p\_value\_lat }\SpecialCharTok{\textless{}} \FloatTok{0.05}\NormalTok{) \{}
    \FunctionTok{print}\NormalTok{(}\StringTok{"Latitude significantly affects area\_diff."}\NormalTok{)}
\NormalTok{\} }\ControlFlowTok{else}\NormalTok{ \{}
    \FunctionTok{print}\NormalTok{(}\StringTok{"Latitude does not significantly affect area\_diff."}\NormalTok{)}
\NormalTok{\}}
\end{Highlighting}
\end{Shaded}

\begin{verbatim}
[1] "Latitude does not significantly affect area_diff."
\end{verbatim}

\begin{Shaded}
\begin{Highlighting}[]
\CommentTok{\# Purpose: Performs a linear regression to examine the effect of latitude on area difference, extracting relevant statistics and checking significance.}

\CommentTok{\# Linear regression with Longitude}
\NormalTok{model\_long }\OtherTok{\textless{}{-}} \FunctionTok{lm}\NormalTok{(area\_diff }\SpecialCharTok{\textasciitilde{}}\NormalTok{ Long }\SpecialCharTok{+}\NormalTok{ genotype\_id, }\AttributeTok{data =}\NormalTok{ coordinates\_merged\_data)}
\NormalTok{summary\_long }\OtherTok{\textless{}{-}} \FunctionTok{summary}\NormalTok{(model\_long)}

\CommentTok{\# Extract coefficients and R{-}squared}
\NormalTok{intercept\_long }\OtherTok{\textless{}{-}} \FunctionTok{coef}\NormalTok{(summary\_long)[}\DecValTok{1}\NormalTok{, }\DecValTok{1}\NormalTok{]}
\NormalTok{slope\_long }\OtherTok{\textless{}{-}} \FunctionTok{coef}\NormalTok{(summary\_long)[}\DecValTok{2}\NormalTok{, }\DecValTok{1}\NormalTok{]}
\NormalTok{r\_squared\_long }\OtherTok{\textless{}{-}}\NormalTok{ summary\_long}\SpecialCharTok{$}\NormalTok{r.squared}
\NormalTok{p\_value\_long }\OtherTok{\textless{}{-}} \FunctionTok{coef}\NormalTok{(summary\_long)[}\DecValTok{2}\NormalTok{, }\DecValTok{4}\NormalTok{]}
\NormalTok{p\_value\_f\_long }\OtherTok{\textless{}{-}}\NormalTok{ summary\_long}\SpecialCharTok{$}\NormalTok{fstatistic[}\DecValTok{1}\NormalTok{] }\CommentTok{\# F{-}statistic value}
\FunctionTok{print}\NormalTok{(summary\_long)}
\end{Highlighting}
\end{Shaded}

\begin{verbatim}

Call:
lm(formula = area_diff ~ Long + genotype_id, data = coordinates_merged_data)

Residuals:
   Min     1Q Median     3Q    Max 
-29382  -4218    614   4153  48129 

Coefficients:
             Estimate Std. Error t value Pr(>|t|)
(Intercept) -499.8901  1509.2140  -0.331    0.741
Long          18.4793    14.2306   1.299    0.195
genotype_id   -0.1080     0.1871  -0.577    0.564

Residual standard error: 8737 on 405 degrees of freedom
Multiple R-squared:  0.004222,  Adjusted R-squared:  -0.0006958 
F-statistic: 0.8585 on 2 and 405 DF,  p-value: 0.4246
\end{verbatim}

\begin{Shaded}
\begin{Highlighting}[]
\ControlFlowTok{if}\NormalTok{ (p\_value\_lat }\SpecialCharTok{\textless{}} \FloatTok{0.05}\NormalTok{) \{}
    \FunctionTok{print}\NormalTok{(}\StringTok{"Longitute significantly affects area\_diff."}\NormalTok{)}
\NormalTok{\} }\ControlFlowTok{else}\NormalTok{ \{}
    \FunctionTok{print}\NormalTok{(}\StringTok{"Longitue does not significantly affect area\_diff."}\NormalTok{)}
\NormalTok{\}}
\end{Highlighting}
\end{Shaded}

\begin{verbatim}
[1] "Longitue does not significantly affect area_diff."
\end{verbatim}

\begin{Shaded}
\begin{Highlighting}[]
\CommentTok{\# Linear regression with Latitude}
\NormalTok{model\_lat }\OtherTok{\textless{}{-}} \FunctionTok{lm}\NormalTok{(area\_diff }\SpecialCharTok{\textasciitilde{}}\NormalTok{ Lat, }\AttributeTok{data =}\NormalTok{ coordinates\_merged\_data)}
\NormalTok{summary\_lat }\OtherTok{\textless{}{-}} \FunctionTok{summary}\NormalTok{(model\_lat)}

\CommentTok{\# Extract coefficients and R{-}squared for Latitude}
\NormalTok{lat\_results }\OtherTok{\textless{}{-}} \FunctionTok{data.frame}\NormalTok{(}
  \AttributeTok{Variable =} \StringTok{"Latitude"}\NormalTok{,}
  \AttributeTok{Intercept =} \FunctionTok{coef}\NormalTok{(summary\_lat)[}\DecValTok{1}\NormalTok{, }\DecValTok{1}\NormalTok{],}
  \AttributeTok{Slope =} \FunctionTok{coef}\NormalTok{(summary\_lat)[}\DecValTok{2}\NormalTok{, }\DecValTok{1}\NormalTok{],}
  \AttributeTok{R\_squared =}\NormalTok{ summary\_lat}\SpecialCharTok{$}\NormalTok{r.squared,}
  \AttributeTok{p\_value\_slope =} \FunctionTok{coef}\NormalTok{(summary\_lat)[}\DecValTok{2}\NormalTok{, }\DecValTok{4}\NormalTok{],}
  \AttributeTok{F\_statistic =}\NormalTok{ summary\_lat}\SpecialCharTok{$}\NormalTok{fstatistic[}\DecValTok{1}\NormalTok{]}
\NormalTok{)}

\CommentTok{\# Linear regression with Longitude}
\NormalTok{model\_long }\OtherTok{\textless{}{-}} \FunctionTok{lm}\NormalTok{(area\_diff }\SpecialCharTok{\textasciitilde{}}\NormalTok{ Long, }\AttributeTok{data =}\NormalTok{ coordinates\_merged\_data)}
\NormalTok{summary\_long }\OtherTok{\textless{}{-}} \FunctionTok{summary}\NormalTok{(model\_long)}

\CommentTok{\# Extract coefficients and R{-}squared for Longitude}
\NormalTok{long\_results }\OtherTok{\textless{}{-}} \FunctionTok{data.frame}\NormalTok{(}
  \AttributeTok{Variable =} \StringTok{"Longitude"}\NormalTok{,}
  \AttributeTok{Intercept =} \FunctionTok{coef}\NormalTok{(summary\_long)[}\DecValTok{1}\NormalTok{, }\DecValTok{1}\NormalTok{],}
  \AttributeTok{Slope =} \FunctionTok{coef}\NormalTok{(summary\_long)[}\DecValTok{2}\NormalTok{, }\DecValTok{1}\NormalTok{],}
  \AttributeTok{R\_squared =}\NormalTok{ summary\_long}\SpecialCharTok{$}\NormalTok{r.squared,}
  \AttributeTok{p\_value\_slope =} \FunctionTok{coef}\NormalTok{(summary\_long)[}\DecValTok{2}\NormalTok{, }\DecValTok{4}\NormalTok{],}
  \AttributeTok{F\_statistic =}\NormalTok{ summary\_long}\SpecialCharTok{$}\NormalTok{fstatistic[}\DecValTok{1}\NormalTok{]}
\NormalTok{)}

\CommentTok{\# Combine the results for Latitude and Longitude}
\NormalTok{regression\_results }\OtherTok{\textless{}{-}} \FunctionTok{rbind}\NormalTok{(lat\_results, long\_results)}

\CommentTok{\# Display the results as a neat table}
\FunctionTok{kable}\NormalTok{(regression\_results, }\AttributeTok{caption =} \StringTok{"Linear Regression Results for Latitude and Longitude"}\NormalTok{, }\AttributeTok{digits =} \DecValTok{4}\NormalTok{)}
\end{Highlighting}
\end{Shaded}

\begin{longtable}[]{@{}
  >{\raggedright\arraybackslash}p{(\linewidth - 12\tabcolsep) * \real{0.0986}}
  >{\raggedright\arraybackslash}p{(\linewidth - 12\tabcolsep) * \real{0.1408}}
  >{\raggedleft\arraybackslash}p{(\linewidth - 12\tabcolsep) * \real{0.1408}}
  >{\raggedleft\arraybackslash}p{(\linewidth - 12\tabcolsep) * \real{0.1127}}
  >{\raggedleft\arraybackslash}p{(\linewidth - 12\tabcolsep) * \real{0.1408}}
  >{\raggedleft\arraybackslash}p{(\linewidth - 12\tabcolsep) * \real{0.1972}}
  >{\raggedleft\arraybackslash}p{(\linewidth - 12\tabcolsep) * \real{0.1690}}@{}}
\caption{Linear Regression Results for Latitude and
Longitude}\tabularnewline
\toprule\noalign{}
\begin{minipage}[b]{\linewidth}\raggedright
\end{minipage} & \begin{minipage}[b]{\linewidth}\raggedright
Variable
\end{minipage} & \begin{minipage}[b]{\linewidth}\raggedleft
Intercept
\end{minipage} & \begin{minipage}[b]{\linewidth}\raggedleft
Slope
\end{minipage} & \begin{minipage}[b]{\linewidth}\raggedleft
R\_squared
\end{minipage} & \begin{minipage}[b]{\linewidth}\raggedleft
p\_value\_slope
\end{minipage} & \begin{minipage}[b]{\linewidth}\raggedleft
F\_statistic
\end{minipage} \\
\midrule\noalign{}
\endfirsthead
\toprule\noalign{}
\begin{minipage}[b]{\linewidth}\raggedright
\end{minipage} & \begin{minipage}[b]{\linewidth}\raggedright
Variable
\end{minipage} & \begin{minipage}[b]{\linewidth}\raggedleft
Intercept
\end{minipage} & \begin{minipage}[b]{\linewidth}\raggedleft
Slope
\end{minipage} & \begin{minipage}[b]{\linewidth}\raggedleft
R\_squared
\end{minipage} & \begin{minipage}[b]{\linewidth}\raggedleft
p\_value\_slope
\end{minipage} & \begin{minipage}[b]{\linewidth}\raggedleft
F\_statistic
\end{minipage} \\
\midrule\noalign{}
\endhead
\bottomrule\noalign{}
\endlastfoot
value & Latitude & -1331.343 & 1.8718 & 0.0000 & 0.9751 & 0.001 \\
value1 & Longitude & -1333.361 & 15.8797 & 0.0034 & 0.2398 & 1.386 \\
\end{longtable}

\subsubsection{Scatter Plots with Regression
Statistics}\label{scatter-plots-with-regression-statistics}

\textbf{Purpose:} Creates scatter plots of latitude and longitude
against area difference, including regression lines and displaying R²
and p-values.

\begin{Shaded}
\begin{Highlighting}[]
\CommentTok{\# Compute linear regression model}
\NormalTok{lm\_lat }\OtherTok{\textless{}{-}} \FunctionTok{lm}\NormalTok{(area\_diff }\SpecialCharTok{\textasciitilde{}}\NormalTok{ Lat, }\AttributeTok{data =}\NormalTok{ coordinates\_merged\_data, }\AttributeTok{na.action =}\NormalTok{ na.omit)}

\CommentTok{\# Extract R² and p{-}value}
\NormalTok{r\_squared\_lat }\OtherTok{\textless{}{-}} \FunctionTok{summary}\NormalTok{(lm\_lat)}\SpecialCharTok{$}\NormalTok{r.squared}
\NormalTok{p\_value\_lat }\OtherTok{\textless{}{-}} \FunctionTok{summary}\NormalTok{(lm\_lat)}\SpecialCharTok{$}\NormalTok{fstatistic[}\DecValTok{1}\NormalTok{] }\SpecialCharTok{\%\textgreater{}\%} 
  \FunctionTok{pf}\NormalTok{(}\AttributeTok{df1 =} \DecValTok{1}\NormalTok{, }\AttributeTok{df2 =}\NormalTok{ lm\_lat}\SpecialCharTok{$}\NormalTok{df.residual, }\AttributeTok{lower.tail =} \ConstantTok{FALSE}\NormalTok{)}

\CommentTok{\# Define y{-}axis limits dynamically (prevents cut{-}off issues)}
\NormalTok{y\_limits }\OtherTok{\textless{}{-}} \FunctionTok{c}\NormalTok{(}\FunctionTok{min}\NormalTok{(coordinates\_merged\_data}\SpecialCharTok{$}\NormalTok{area\_diff, }\AttributeTok{na.rm =} \ConstantTok{TRUE}\NormalTok{) }\SpecialCharTok{*} \FloatTok{1.1}\NormalTok{, }
              \FunctionTok{max}\NormalTok{(coordinates\_merged\_data}\SpecialCharTok{$}\NormalTok{area\_diff, }\AttributeTok{na.rm =} \ConstantTok{TRUE}\NormalTok{) }\SpecialCharTok{*} \FloatTok{1.1}\NormalTok{)}

\CommentTok{\# Create the scatter plot}
\NormalTok{corr\_plot\_lat }\OtherTok{\textless{}{-}} \FunctionTok{ggplot}\NormalTok{(coordinates\_merged\_data, }
                        \FunctionTok{aes}\NormalTok{(}\AttributeTok{x =}\NormalTok{ Lat, }\AttributeTok{y =}\NormalTok{ area\_diff, }\AttributeTok{color =}\NormalTok{ preference)) }\SpecialCharTok{+}
  
  \CommentTok{\# Scatter points}
  \FunctionTok{geom\_point}\NormalTok{(}\AttributeTok{size =} \DecValTok{1}\NormalTok{, }\AttributeTok{alpha =} \FloatTok{0.8}\NormalTok{) }\SpecialCharTok{+}
  
  \CommentTok{\# Custom color scale}
  \FunctionTok{scale\_colour\_manual}\NormalTok{(}
    \AttributeTok{values =} \FunctionTok{c}\NormalTok{(}\StringTok{"orange2"}\NormalTok{, }\StringTok{"\#0072B2"}\NormalTok{),}
    \AttributeTok{name =} \StringTok{"BX feedback"}\NormalTok{,}
    \AttributeTok{labels =} \FunctionTok{c}\NormalTok{(}\StringTok{"negative"}\NormalTok{, }\StringTok{"positive"}\NormalTok{)}
\NormalTok{  ) }\SpecialCharTok{+}
  
  \CommentTok{\# Regression line}
  \FunctionTok{geom\_smooth}\NormalTok{(}\AttributeTok{method =} \StringTok{"lm"}\NormalTok{, }\AttributeTok{se =} \ConstantTok{FALSE}\NormalTok{, }\AttributeTok{linetype =} \StringTok{"solid"}\NormalTok{, }
              \AttributeTok{color =} \StringTok{"red4"}\NormalTok{, }\AttributeTok{linewidth =} \DecValTok{1}\NormalTok{) }\SpecialCharTok{+}
  
  \CommentTok{\# Labels}
  \FunctionTok{labs}\NormalTok{(}
    \AttributeTok{title =} \StringTok{"Phenotype vs. Latitude with Regression Line"}\NormalTok{,}
    \AttributeTok{x =} \StringTok{"Latitude"}\NormalTok{,}
    \AttributeTok{y =} \StringTok{"Phenotype (Area Diff)"}
\NormalTok{  ) }\SpecialCharTok{+}
  
  \CommentTok{\# Theme modifications for A4 compatibility}
  \FunctionTok{theme\_classic}\NormalTok{() }\SpecialCharTok{+}
  \FunctionTok{ylim}\NormalTok{(y\_limits) }\SpecialCharTok{+}
  \FunctionTok{theme}\NormalTok{(}
    \AttributeTok{legend.position =} \FunctionTok{c}\NormalTok{(}\FloatTok{0.97}\NormalTok{, }\FloatTok{0.03}\NormalTok{),  }\CommentTok{\# Bottom{-}right inside plot}
    \AttributeTok{legend.justification =} \FunctionTok{c}\NormalTok{(}\DecValTok{1}\NormalTok{, }\DecValTok{0}\NormalTok{),  }
    \AttributeTok{legend.background =} \FunctionTok{element\_rect}\NormalTok{(}\AttributeTok{fill =} \StringTok{"white"}\NormalTok{, }\AttributeTok{color =} \StringTok{"black"}\NormalTok{),}
    \AttributeTok{legend.title =} \FunctionTok{element\_text}\NormalTok{(}\AttributeTok{size =} \DecValTok{6}\NormalTok{),}
    \AttributeTok{legend.text =} \FunctionTok{element\_text}\NormalTok{(}\AttributeTok{size =} \DecValTok{5}\NormalTok{),}
    \AttributeTok{panel.border =} \FunctionTok{element\_rect}\NormalTok{(}\AttributeTok{color =} \StringTok{"black"}\NormalTok{, }\AttributeTok{fill =} \ConstantTok{NA}\NormalTok{, }\AttributeTok{size =} \FloatTok{0.5}\NormalTok{),}
    \AttributeTok{plot.title =} \FunctionTok{element\_text}\NormalTok{(}\AttributeTok{size =} \DecValTok{8}\NormalTok{, }\AttributeTok{face =} \StringTok{"bold"}\NormalTok{),}
    \AttributeTok{axis.title =} \FunctionTok{element\_text}\NormalTok{(}\AttributeTok{size =} \DecValTok{7}\NormalTok{),}
    \AttributeTok{axis.text =} \FunctionTok{element\_text}\NormalTok{(}\AttributeTok{size =} \DecValTok{6}\NormalTok{)}
\NormalTok{  ) }\SpecialCharTok{+}
  
  \CommentTok{\# Add R² and p{-}value annotation}
  \FunctionTok{annotate}\NormalTok{(}
    \StringTok{"text"}\NormalTok{, }\AttributeTok{x =} \FunctionTok{max}\NormalTok{(coordinates\_merged\_data}\SpecialCharTok{$}\NormalTok{Lat, }\AttributeTok{na.rm =} \ConstantTok{TRUE}\NormalTok{), }
    \AttributeTok{y =} \FunctionTok{min}\NormalTok{(y\_limits) }\SpecialCharTok{+}\NormalTok{ (}\FunctionTok{max}\NormalTok{(y\_limits) }\SpecialCharTok{{-}} \FunctionTok{min}\NormalTok{(y\_limits)) }\SpecialCharTok{*} \FloatTok{0.85}\NormalTok{,  }\CommentTok{\# Position dynamically}
    \AttributeTok{label =} \FunctionTok{sprintf}\NormalTok{(}\StringTok{"R² = \%.2f}\SpecialCharTok{\textbackslash{}n}\StringTok{p (F) = \%.3g"}\NormalTok{, r\_squared\_lat, p\_value\_lat),}
    \AttributeTok{hjust =} \DecValTok{1}\NormalTok{, }\AttributeTok{vjust =} \DecValTok{0}\NormalTok{, }\AttributeTok{size =} \DecValTok{3}\NormalTok{, }\AttributeTok{color =} \StringTok{"black"}
\NormalTok{  )}

\CommentTok{\# Display the plot}
\FunctionTok{print}\NormalTok{(corr\_plot\_lat)}
\end{Highlighting}
\end{Shaded}

\begin{Shaded}
\begin{Highlighting}[]
\CommentTok{\# Compute linear regression model}
\NormalTok{lm\_long }\OtherTok{\textless{}{-}} \FunctionTok{lm}\NormalTok{(area\_diff }\SpecialCharTok{\textasciitilde{}}\NormalTok{ Long, }\AttributeTok{data =}\NormalTok{ coordinates\_merged\_data, }
              \AttributeTok{na.action =}\NormalTok{ na.omit)}

\CommentTok{\# Extract R² and p{-}value}
\NormalTok{r\_squared\_long }\OtherTok{\textless{}{-}} \FunctionTok{summary}\NormalTok{(lm\_long)}\SpecialCharTok{$}\NormalTok{r.squared}
\NormalTok{p\_value\_long }\OtherTok{\textless{}{-}} \FunctionTok{summary}\NormalTok{(lm\_long)}\SpecialCharTok{$}\NormalTok{fstatistic[}\DecValTok{1}\NormalTok{] }\SpecialCharTok{\%\textgreater{}\%} 
  \FunctionTok{pf}\NormalTok{(}\AttributeTok{df1 =} \DecValTok{1}\NormalTok{, }\AttributeTok{df2 =}\NormalTok{ lm\_long}\SpecialCharTok{$}\NormalTok{df.residual, }\AttributeTok{lower.tail =} \ConstantTok{FALSE}\NormalTok{)}

\CommentTok{\# Define y{-}axis limits dynamically (prevents cut{-}off issues)}
\NormalTok{y\_limits }\OtherTok{\textless{}{-}} \FunctionTok{c}\NormalTok{(}\FunctionTok{min}\NormalTok{(coordinates\_merged\_data}\SpecialCharTok{$}\NormalTok{area\_diff, }\AttributeTok{na.rm =} \ConstantTok{TRUE}\NormalTok{) }\SpecialCharTok{*} \FloatTok{1.1}\NormalTok{, }
              \FunctionTok{max}\NormalTok{(coordinates\_merged\_data}\SpecialCharTok{$}\NormalTok{area\_diff, }\AttributeTok{na.rm =} \ConstantTok{TRUE}\NormalTok{) }\SpecialCharTok{*} \FloatTok{1.1}\NormalTok{)}

\CommentTok{\# Create the scatter plot}
\NormalTok{corr\_plot\_long }\OtherTok{\textless{}{-}} \FunctionTok{ggplot}\NormalTok{(coordinates\_merged\_data, }\FunctionTok{aes}\NormalTok{(}\AttributeTok{x =}\NormalTok{ Long, }\AttributeTok{y =}\NormalTok{ area\_diff, }
                                                      \AttributeTok{color =}\NormalTok{ preference)) }\SpecialCharTok{+}
  
  \CommentTok{\# Scatter points}
  \FunctionTok{geom\_point}\NormalTok{(}\AttributeTok{size =} \DecValTok{1}\NormalTok{, }\AttributeTok{alpha =} \FloatTok{0.8}\NormalTok{) }\SpecialCharTok{+}
  
  \CommentTok{\# Custom color scale}
  \FunctionTok{scale\_colour\_manual}\NormalTok{(}
    \AttributeTok{values =} \FunctionTok{c}\NormalTok{(}\StringTok{"orange2"}\NormalTok{, }\StringTok{"\#0072B2"}\NormalTok{),}
    \AttributeTok{name =} \StringTok{"BX feedback"}\NormalTok{,}
    \AttributeTok{labels =} \FunctionTok{c}\NormalTok{(}\StringTok{"negative"}\NormalTok{, }\StringTok{"positive"}\NormalTok{)}
\NormalTok{  ) }\SpecialCharTok{+}
  
  \CommentTok{\# Regression line}
  \FunctionTok{geom\_smooth}\NormalTok{(}\AttributeTok{method =} \StringTok{"lm"}\NormalTok{, }\AttributeTok{se =} \ConstantTok{FALSE}\NormalTok{, }\AttributeTok{linetype =} \StringTok{"solid"}\NormalTok{, }
              \AttributeTok{color =} \StringTok{"red4"}\NormalTok{, }\AttributeTok{linewidth =} \DecValTok{1}\NormalTok{) }\SpecialCharTok{+}
  
  \CommentTok{\# Labels}
  \FunctionTok{labs}\NormalTok{(}
    \AttributeTok{title =} \StringTok{"Phenotype vs. Longitude with Regression Line"}\NormalTok{,}
    \AttributeTok{x =} \StringTok{"Longitude"}\NormalTok{,}
    \AttributeTok{y =} \StringTok{"Phenotype (Area Diff)"}
\NormalTok{  ) }\SpecialCharTok{+}
  
  \CommentTok{\# Theme modifications for A4 compatibility}
  \FunctionTok{theme\_classic}\NormalTok{() }\SpecialCharTok{+}
  \FunctionTok{ylim}\NormalTok{(y\_limits) }\SpecialCharTok{+}
  \FunctionTok{theme}\NormalTok{(}
    \AttributeTok{legend.position =} \FunctionTok{c}\NormalTok{(}\FloatTok{0.97}\NormalTok{, }\FloatTok{0.03}\NormalTok{),  }\CommentTok{\# Bottom{-}right inside plot (same as lat plot)}
    \AttributeTok{legend.justification =} \FunctionTok{c}\NormalTok{(}\DecValTok{1}\NormalTok{, }\DecValTok{0}\NormalTok{),  }
    \AttributeTok{legend.background =} \FunctionTok{element\_rect}\NormalTok{(}\AttributeTok{fill =} \StringTok{"white"}\NormalTok{, }\AttributeTok{color =} \StringTok{"black"}\NormalTok{),}
    \AttributeTok{legend.title =} \FunctionTok{element\_text}\NormalTok{(}\AttributeTok{size =} \DecValTok{6}\NormalTok{),}
    \AttributeTok{legend.text =} \FunctionTok{element\_text}\NormalTok{(}\AttributeTok{size =} \DecValTok{5}\NormalTok{),}
    \AttributeTok{panel.border =} \FunctionTok{element\_rect}\NormalTok{(}\AttributeTok{color =} \StringTok{"black"}\NormalTok{, }\AttributeTok{fill =} \ConstantTok{NA}\NormalTok{, }\AttributeTok{size =} \FloatTok{0.5}\NormalTok{),}
    \AttributeTok{plot.title =} \FunctionTok{element\_text}\NormalTok{(}\AttributeTok{size =} \DecValTok{8}\NormalTok{, }\AttributeTok{face =} \StringTok{"bold"}\NormalTok{),}
    \AttributeTok{axis.title =} \FunctionTok{element\_text}\NormalTok{(}\AttributeTok{size =} \DecValTok{7}\NormalTok{),}
    \AttributeTok{axis.text =} \FunctionTok{element\_text}\NormalTok{(}\AttributeTok{size =} \DecValTok{6}\NormalTok{)}
\NormalTok{  ) }\SpecialCharTok{+}
  
  \CommentTok{\# Add R² and p{-}value annotation}
  \FunctionTok{annotate}\NormalTok{(}
    \StringTok{"text"}\NormalTok{, }\AttributeTok{x =} \FunctionTok{max}\NormalTok{(coordinates\_merged\_data}\SpecialCharTok{$}\NormalTok{Long, }\AttributeTok{na.rm =} \ConstantTok{TRUE}\NormalTok{), }
    \AttributeTok{y =} \FunctionTok{min}\NormalTok{(y\_limits) }\SpecialCharTok{+}\NormalTok{ (}\FunctionTok{max}\NormalTok{(y\_limits) }\SpecialCharTok{{-}} \FunctionTok{min}\NormalTok{(y\_limits)) }\SpecialCharTok{*} \FloatTok{0.85}\NormalTok{,  }\CommentTok{\# Position dynamically}
    \AttributeTok{label =} \FunctionTok{sprintf}\NormalTok{(}\StringTok{"R² = \%.2f}\SpecialCharTok{\textbackslash{}n}\StringTok{p (F) = \%.3g"}\NormalTok{, r\_squared\_long, p\_value\_long),}
    \AttributeTok{hjust =} \DecValTok{1}\NormalTok{, }\AttributeTok{vjust =} \DecValTok{0}\NormalTok{, }\AttributeTok{size =} \DecValTok{3}\NormalTok{, }\AttributeTok{color =} \StringTok{"black"}
\NormalTok{  )}

\CommentTok{\# Display the plot}
\FunctionTok{print}\NormalTok{(corr\_plot\_long)}
\end{Highlighting}
\end{Shaded}

\begin{figure}[H]

{\centering \pandocbounded{\includegraphics[keepaspectratio]{01_Accession_phenotypes_files/figure-pdf/FigS1a-Corr-Plots-1.pdf}}

}

\caption{Correlation of geographic distribution of accessions used in
the GWASand their growth feedbacks.}

\end{figure}%

\subsection{Plot Accession Locations on a World
Map}\label{plot-accession-locations-on-a-world-map}

\textbf{Purpos}e: Prepares the world map data and creates a base plot
for plotting the locations of positive and negative accessions.

\begin{Shaded}
\begin{Highlighting}[]
\CommentTok{\# Load necessary libraries}
\FunctionTok{library}\NormalTok{(maps)}
\CommentTok{\# Purpose: Loads the \textasciigrave{}maps\textasciigrave{} package for map plotting.}

\CommentTok{\# Obtain world map data}
\NormalTok{world\_map }\OtherTok{\textless{}{-}} \FunctionTok{map\_data}\NormalTok{(}\StringTok{"world"}\NormalTok{)}
\CommentTok{\# Purpose: Retrieves the data necessary to create a world map.}

\CommentTok{\# Subset the world map data for the region of interest}
\NormalTok{World\_map\_small }\OtherTok{\textless{}{-}} \FunctionTok{subset}\NormalTok{(world\_map, long }\SpecialCharTok{\textgreater{}} \SpecialCharTok{{-}}\DecValTok{130} \SpecialCharTok{\&}\NormalTok{ long }\SpecialCharTok{\textless{}} \DecValTok{90}\NormalTok{)}
\NormalTok{World\_map\_small }\OtherTok{\textless{}{-}} \FunctionTok{subset}\NormalTok{(World\_map\_small, lat }\SpecialCharTok{\textgreater{}} \SpecialCharTok{{-}}\DecValTok{0} \SpecialCharTok{\&}\NormalTok{ lat }\SpecialCharTok{\textless{}} \DecValTok{110}\NormalTok{)}
\CommentTok{\# Purpose: Limits the map data to the specific region of interest by filtering longitude and latitude ranges.}

\CommentTok{\# Create a base plot with ggplot2}
\NormalTok{p }\OtherTok{\textless{}{-}} \FunctionTok{ggplot}\NormalTok{() }\SpecialCharTok{+} \FunctionTok{coord\_fixed}\NormalTok{() }\SpecialCharTok{+}
  \FunctionTok{xlab}\NormalTok{(}\StringTok{"Longitude"}\NormalTok{) }\SpecialCharTok{+} \FunctionTok{ylab}\NormalTok{(}\StringTok{"Latitude"}\NormalTok{)}
\CommentTok{\# Purpose: Sets up the basic structure for the map plot with fixed coordinates.}

\CommentTok{\# Add map to base plot}
\NormalTok{base\_world\_messy }\OtherTok{\textless{}{-}}\NormalTok{ p }\SpecialCharTok{+} \FunctionTok{geom\_polygon}\NormalTok{(}\AttributeTok{data =}\NormalTok{ World\_map\_small, }\FunctionTok{aes}\NormalTok{(}\AttributeTok{x =}\NormalTok{ long, }
                                                                 \AttributeTok{y =}\NormalTok{ lat, }
                                                                 \AttributeTok{group =}\NormalTok{ group), }
                                     \AttributeTok{colour =} \StringTok{"grey"}\NormalTok{, }\AttributeTok{fill =} \StringTok{"grey"}\NormalTok{)}
\CommentTok{\# Purpose: Adds the filtered map data to the plot as polygons representing countries.}

\CommentTok{\# Display the messy base world map}
\CommentTok{\#print(base\_world\_messy)}
\CommentTok{\# Purpose: Displays the initial, unrefined world map.}

\CommentTok{\# Strip the map down to make it clean and beautiful}
\NormalTok{cleanup }\OtherTok{\textless{}{-}} \FunctionTok{theme}\NormalTok{(}\AttributeTok{panel.grid.major =} \FunctionTok{element\_blank}\NormalTok{(), }
                 \AttributeTok{panel.grid.minor =} \FunctionTok{element\_blank}\NormalTok{(), }
                 \AttributeTok{panel.background =} \FunctionTok{element\_rect}\NormalTok{(}\AttributeTok{fill =} \StringTok{\textquotesingle{}white\textquotesingle{}}\NormalTok{, }
                                                 \AttributeTok{colour =} \StringTok{\textquotesingle{}white\textquotesingle{}}\NormalTok{), }
                 \AttributeTok{axis.line =} \FunctionTok{element\_line}\NormalTok{(}\AttributeTok{colour =} \StringTok{"grey50"}\NormalTok{))}
\CommentTok{\# Purpose: Cleans up the map by removing unnecessary elements and setting a minimalist theme.}

\NormalTok{base\_world }\OtherTok{\textless{}{-}}\NormalTok{ base\_world\_messy }\SpecialCharTok{+}\NormalTok{ cleanup}
\CommentTok{\# Purpose: Applies the cleaned{-}up theme to the world map.}

\CommentTok{\# Display the clean base world map}
\CommentTok{\#print(base\_world)}
\CommentTok{\# Purpose: Displays the refined world map.}

\CommentTok{\# Separate data for layers}
\NormalTok{df\_layer\_1 }\OtherTok{\textless{}{-}}\NormalTok{ coordinates\_merged\_data[coordinates\_merged\_data}\SpecialCharTok{$}\NormalTok{preference }\SpecialCharTok{==} \StringTok{"plus"}\NormalTok{, ]}
\NormalTok{df\_layer\_2 }\OtherTok{\textless{}{-}}\NormalTok{ coordinates\_merged\_data[coordinates\_merged\_data}\SpecialCharTok{$}\NormalTok{preference }\SpecialCharTok{==} \StringTok{"minus"}\NormalTok{, ]}
\CommentTok{\# Purpose: Separates the coordinates data into two layers based on whether the response is positive or negative.}
\end{Highlighting}
\end{Shaded}

Plot the map and overlay with the accessions used in the GWAS

\begin{Shaded}
\begin{Highlighting}[]
\CommentTok{\# Add simple data points to the map}
\NormalTok{ACC\_map }\OtherTok{\textless{}{-}}\NormalTok{ base\_world }\SpecialCharTok{+}
  \FunctionTok{geom\_point}\NormalTok{(}
    \AttributeTok{data =}\NormalTok{ df\_layer\_2,}
    \FunctionTok{aes}\NormalTok{(}\AttributeTok{x =}\NormalTok{ Long, }\AttributeTok{y =}\NormalTok{ Lat, }\AttributeTok{fill =} \StringTok{"minus"}\NormalTok{),}
    \AttributeTok{colour =} \StringTok{"black"}\NormalTok{,}
    \AttributeTok{pch =} \DecValTok{21}\NormalTok{,}
    \AttributeTok{size =} \DecValTok{2}\NormalTok{,}
    \AttributeTok{alpha =} \FloatTok{0.5}
\NormalTok{  ) }\SpecialCharTok{+}
  \FunctionTok{geom\_point}\NormalTok{(}
    \AttributeTok{data =}\NormalTok{ df\_layer\_1,}
    \FunctionTok{aes}\NormalTok{(}\AttributeTok{x =}\NormalTok{ Long, }\AttributeTok{y =}\NormalTok{ Lat, }\AttributeTok{fill =} \StringTok{"plus"}\NormalTok{),}
    \AttributeTok{colour =} \StringTok{"black"}\NormalTok{,}
    \AttributeTok{pch =} \DecValTok{21}\NormalTok{,}
    \AttributeTok{size =} \DecValTok{2}\NormalTok{,}
    \AttributeTok{alpha =} \FloatTok{0.5}
\NormalTok{  ) }\SpecialCharTok{+}
  \FunctionTok{scale\_fill\_manual}\NormalTok{(}
    \AttributeTok{values =} \FunctionTok{c}\NormalTok{(}\StringTok{"plus"} \OtherTok{=} \StringTok{"\#0072B2"}\NormalTok{, }\StringTok{"minus"} \OtherTok{=} \StringTok{"orange2"}\NormalTok{),}
    \AttributeTok{name =} \StringTok{"Growth feedback"}\NormalTok{,}
    \AttributeTok{labels =} \FunctionTok{c}\NormalTok{(}\StringTok{"negative"}\NormalTok{, }\StringTok{"positive"}\NormalTok{)}
\NormalTok{  ) }\SpecialCharTok{+}
  \FunctionTok{theme}\NormalTok{(}
    \AttributeTok{legend.position =} \FunctionTok{c}\NormalTok{(}\FloatTok{0.01}\NormalTok{, }\FloatTok{0.01}\NormalTok{),  }\CommentTok{\# Bottom{-}left corner}
    \AttributeTok{legend.justification =} \FunctionTok{c}\NormalTok{(}\DecValTok{0}\NormalTok{, }\DecValTok{0}\NormalTok{),  }\CommentTok{\# Align legend to bottom{-}left}
    \AttributeTok{legend.background =} \FunctionTok{element\_rect}\NormalTok{(}\AttributeTok{fill =} \StringTok{"white"}\NormalTok{, }\AttributeTok{color =} \StringTok{"black"}\NormalTok{),  }\CommentTok{\# Optional: Add legend background}
    \AttributeTok{legend.text =} \FunctionTok{element\_text}\NormalTok{(}\AttributeTok{size =} \DecValTok{7}\NormalTok{),  }\CommentTok{\# Adjust legend text size}
    \AttributeTok{legend.title =} \FunctionTok{element\_text}\NormalTok{(}\AttributeTok{size =} \DecValTok{7}\NormalTok{, }\AttributeTok{face =} \StringTok{"bold"}\NormalTok{)  }\CommentTok{\# Adjust legend title}
\NormalTok{  )}

\CommentTok{\# Display the map with data points}
\FunctionTok{print}\NormalTok{(ACC\_map)}
\end{Highlighting}
\end{Shaded}

\begin{figure}[H]

{\centering \pandocbounded{\includegraphics[keepaspectratio]{01_Accession_phenotypes_files/figure-pdf/Fig1d-ACC-Geographic-1.pdf}}

}

\caption{Geographic distribution of accessions used in the GWAS.}

\end{figure}%

\begin{Shaded}
\begin{Highlighting}[]
\CommentTok{\# Purpose: Displays the world map with the positive and negative accession data points.}

\CommentTok{\# Save the plot as a PDF}
\FunctionTok{ggsave}\NormalTok{(}\StringTok{"ACC\_response\_map\_white.pdf"}\NormalTok{, }\AttributeTok{plot =}\NormalTok{ ACC\_map, }\AttributeTok{height =} \DecValTok{15}\NormalTok{, }\AttributeTok{width =} \DecValTok{20}\NormalTok{, }\AttributeTok{units =} \StringTok{"cm"}\NormalTok{, }\AttributeTok{path =} \StringTok{"Output/Figures/"}\NormalTok{)}
\CommentTok{\# Purpose: Saves the map with data points as a PDF file.}
\end{Highlighting}
\end{Shaded}

\section{Session information}\label{session-information}

\textbf{Purpose}: Print the session information, including details about
the R version, loaded packages, and the system environment. This is
useful for reproducibility, as it documents the computational
environment in which the analysis was performed.

\begin{Shaded}
\begin{Highlighting}[]
\CommentTok{\# Get session information}
\NormalTok{SI }\OtherTok{\textless{}{-}} \FunctionTok{sessionInfo}\NormalTok{()}

\CommentTok{\# Extract the attached packages and their versions}
\NormalTok{attached\_packages }\OtherTok{\textless{}{-}}\NormalTok{ SI}\SpecialCharTok{$}\NormalTok{otherPkgs}

\CommentTok{\# Print the R version}
\FunctionTok{cat}\NormalTok{(}\StringTok{"R version:"}\NormalTok{, SI}\SpecialCharTok{$}\NormalTok{R.version}\SpecialCharTok{$}\NormalTok{version.string, }\StringTok{"}\SpecialCharTok{\textbackslash{}n\textbackslash{}n}\StringTok{"}\NormalTok{)}
\end{Highlighting}
\end{Shaded}

\begin{verbatim}
R version: R version 4.4.2 (2024-10-31) 
\end{verbatim}

\begin{Shaded}
\begin{Highlighting}[]
\CommentTok{\# Print the attached packages and their versions}
\FunctionTok{cat}\NormalTok{(}\StringTok{"Attached packages:}\SpecialCharTok{\textbackslash{}n}\StringTok{"}\NormalTok{)}
\end{Highlighting}
\end{Shaded}

\begin{verbatim}
Attached packages:
\end{verbatim}

\begin{Shaded}
\begin{Highlighting}[]
\FunctionTok{invisible}\NormalTok{(}\FunctionTok{lapply}\NormalTok{(attached\_packages, }\ControlFlowTok{function}\NormalTok{(pkg) \{}
  \FunctionTok{cat}\NormalTok{(pkg}\SpecialCharTok{$}\NormalTok{Package, }\StringTok{"{-}"}\NormalTok{, pkg}\SpecialCharTok{$}\NormalTok{Version, }\StringTok{"}\SpecialCharTok{\textbackslash{}n}\StringTok{"}\NormalTok{)}
\NormalTok{\}))}
\end{Highlighting}
\end{Shaded}

\begin{verbatim}
maps - 3.4.2.1 
dtw - 1.23-1 
proxy - 0.4-27 
fstcore - 0.9.18 
readxl - 1.4.3 
openxlsx - 4.2.7.1 
circlize - 0.4.16 
ComplexHeatmap - 2.22.0 
ggdendro - 0.2.0 
knitr - 1.49 
cowplot - 1.1.3 
MESS - 0.5.12 
ggpubr - 0.6.0 
DescTools - 0.99.58 
magrittr - 2.0.3 
gwaR - 0.6 
lubridate - 1.9.4 
forcats - 1.0.0 
stringr - 1.5.1 
dplyr - 1.1.4 
purrr - 1.0.2 
readr - 2.1.5 
tidyr - 1.3.1 
tibble - 3.2.1 
ggplot2 - 3.5.1 
tidyverse - 2.0.0 
\end{verbatim}




\end{document}
